\chapter{Modal logic as a semantic perspective on structures}
\label{sec:p08-src-025-501}

This presentation follows Blackburn and van Benthem's institutional manuscript
\emph{Modal Logic: A Semantic Perspective}, ILLC PP-2006-30
\cite{BlackburnVanBenthem2006ILLC}. Its organizing question is what a modal
language can express about a structure, and how that expressivity affects
reasoning. Relational, algebraic and topological interpretations are different
semantic settings. The relational setting begins with an internal viewpoint:
formulas are evaluated at a selected state and explore its accessible states.
That viewpoint supports temporal, epistemic, computational and ordinary
extensional interpretations. Possible worlds supply one interpretation of
the states, not a restriction built into the mathematical definition.

The presentation distinguishes languages, classes of models and classes of
frames. Model checking asks about truth in one given structure; satisfiability
asks whether a suitable structure exists; frame validity quantifies over
valuations. None of these tasks becomes another merely because the same
modal symbols appear. We use ordinary set-based structures and classical
metatheory; uses of arbitrary-cardinal compactness and witness choices below
assume ZFC.

\section{Basic language and relational semantics}
\label{sec:p08-bvb501-basic}
\index{basic modal language}
\index{signature}
\index{similarity type}
\index{domain}
\index{points}
\index{worlds}
\index{valuation}
\index{satisfaction definition}
\index{satisfied}
\index{satisfiable}
\index{true}
\index{valid}
\index{semantic consequence}
\index{globally true}
\index{globally satisfied}
\index{box}
\index{diamond}
\index{modal operator depth}
\index{subformulas}
\index{internal perspective}
\index{possible world semantics}
\index{Kripke model}

\begin{hbdefinition}{def:p08-bvb501-language}{Source language, models and satisfaction}
Fix a set \(\mathrm{PROP}\) of proposition symbols and a set
\(\mathrm{MOD}\) of modality names. Formulas are generated by
\[
 \varphi ::= p\mid\topv\mid\bottom\mid\neg\varphi
 \mid(\varphi\circ\psi)\mid\langle m\rangle\varphi
 \mid[m]\varphi,
 \qquad\circ\in\{\land,\lor,\to,\leftrightarrow\}.
\]
The redundant Boolean primitives preserve the source presentation. For a
single modality, abbreviate its diamond and box by \(\Diamond\) and \(\Box\).
A model is \(M=(W,(R_m)_{m\in\mathrm{MOD}},V)\), with
\(W\ne\varnothing\), \(R_m\subseteq W^2\), and
\(V:\mathrm{PROP}\to\mathcal P(W)\). Its underlying frame forgets
\(V\). Truth at \(w\in W\) uses the classical Boolean clauses and
\[
 \begin{aligned}
 M,w\entails p&\quad\text{iff}\quad w\in V(p),\\
 M,w\entails\langle m\rangle\varphi
  &\quad\text{iff}\quad\exists v\,(wR_mv\ \text{and}\ M,v\entails\varphi),\\
 M,w\entails[m]\varphi
  &\quad\text{iff}\quad\forall v\,(wR_mv\ \text{implies}\ M,v\entails\varphi).
 \end{aligned}
\]
The modal depth \(\operatorname{md}(\varphi)\) is zero on atoms and
Boolean constants, unchanged by negation, the maximum on binary Boolean
combinations, and one plus the argument's depth at a modality.
\end{hbdefinition}

\begin{hbdefinition}{def:p08-bvb501-consequence}{Local and global semantic questions}
Write \(M\entails\varphi\) for truth at every state of \(M\).
A formula is valid when this holds in every model, and satisfiable when it
holds at some state of some model. A set \(\Sigma\) is locally satisfiable
when all its formulas hold at one state. Its local consequence relation is
\(\Sigma\entails\varphi\) iff every pointed model satisfying all of
\(\Sigma\) satisfies \(\varphi\). Global premise truth, by contrast,
requires all premises at every state; it is a different condition.
\end{hbdefinition}

The box at a state without successors is true for every argument, including
\(\bottom\); the diamond is false for every argument there. Quantification
is over successors of the current state. The duality
\([m]\varphi\leftrightarrow\neg\langle m\rangle\neg\varphi\)
follows from these clauses. For polarity calculations expand implication
and equivalence into negation, conjunction and disjunction first: the
antecedent of an implication reverses polarity, even though no negation
is visibly printed in the primitive implication notation.

\subsection{The source's temporal, process, epistemic and extensional examples}
\label{sec:p08-bvb501-graph-examples}
\index{tense logic}
\index{epistemic logic}
\index{description logic}
\index{extensional interpretation}
\index{intensional interpretation}
\index{isomorphism}
\index{Kripke semantics}

\begin{hbexample}{def:p08-bvb501-time}{Five temporal states}
Take \(W=\{t_1,\ldots,t_5\}\), with \(t_iRt_j\) iff \(i<j\),
\(V(p)=\{t_2,t_4\}\) and \(V(q)=\{t_4,t_5\}\).
The drawn consecutive arrows in the source's Figure~1 stand for their
transitive closure. Thus \(\Diamond(p\land q)\) holds at
\(t_1,t_2,t_3\), whereas \(\Diamond q\) holds also at \(t_4\).
Every \(\Box\varphi\) holds at \(t_5\). Prior's future diamond
\(F\) has this forward interpretation; a past diamond \(P\) uses the
converse relation, making \(P(p\land q)\) true at \(t_5\).
The temporal vocabulary locates a claim relative to the point of evaluation.
\end{hbexample}

\begin{hbexample}{def:p08-bvb501-process}{A labelled finite-state process}
The source's Figure~2 has states \(s,u,t\),
\(R_a=\{(s,u),(u,u)\}\) and
\(R_b=\{(u,t),(t,t)\}\), with initial state \(s\) and accepting
state \(t\). For positive integers \(n,m\),
\(\langle a\rangle^n\langle b\rangle^m\topv\) holds at \(s\).
The accepting paths spell exactly \(a^n b^m\), \(n,m>0\).
Modalities here denote symbol-reading transitions. A formula ending in
\(\topv\) asserts a path exists; acceptance must separately specify its
endpoint, for example by an atom true only at \(t\). In particular,
\(\langle a\rangle\topv\) is true at \(s\) even though the word
\(a\) alone is not accepted. Relations may be arbitrary; their being
partial functions in this example is not a standing modal assumption.
\end{hbexample}

\begin{hbexample}{def:p08-bvb501-epistemic}{Knowledge and compatible alternatives}
In the four-state epistemic diagram in the source's Figure~3, name the states
\(e,a,b,d\), take the accessibility relation to be all of \(W^2\), and
let \(q\) hold everywhere, \(p\) only at \(a\), and \(r\) only at
\(d\). At the displayed current state \(e\),
\(\Diamond(p\land q)\), \(\Diamond(r\land q)\) and \(\Box q\)
hold, while \(\Diamond(p\land q\land r)\) does not.
Read the diamond as compatibility with the agent's information and the
box as knowledge. Hintikka's notation \(Kq\) is the latter reading.
The distinction between a state and a proposition true there is preserved:
\(e\) names a state, not a new proposition letter.
\end{hbexample}

\begin{hbexample}{def:p08-bvb501-individuals}{Ordinary people and two relations}
The source's Figure~4 uses Judy, Johnny, Terry and Frank as ordinary
individuals. Its relations are
\[
\begin{aligned}
 R_{\mathrm{LOVES}}&=\{(\mathrm{Judy},\mathrm{Johnny}),
   (\mathrm{Terry},\mathrm{Judy}),(\mathrm{Frank},\mathrm{Judy})\},\\
 R_{\mathrm{DETESTS}}&=\{(\mathrm{Johnny},\mathrm{Judy}),
   (\mathrm{Terry},\mathrm{Frank}),(\mathrm{Frank},\mathrm{Terry})\}.
\end{aligned}
\]
At Terry, \(\langle\mathrm{LOVES}\rangle\topv\land
\langle\mathrm{DETESTS}\rangle\langle\mathrm{LOVES}\rangle\topv\)
is true: Judy witnesses the first conjunct and Frank, then Judy, witness
the second. The graph is extensional information about people; no possible
world interpretation is needed. Description-logic concept languages can use
the same relational pattern for representing such information.
\end{hbexample}

These examples give distinct interpretations of one mathematical semantics.
The manuscript credits Prior for tense logic, Hintikka for knowledge and
belief, and Kripke, Hintikka and Kanger as pioneers of relational modal
semantics. Its historical discussion also mentions Prior and Montague and
refers detailed priority questions to Goldblatt. This attribution is the
source's historical account, not a claim that the displayed elementary
examples reconstruct all those authors' distinct systems
\cite[pp.~5--9]{BlackburnVanBenthem2006ILLC}.

\begin{hbdefinition}{def:p08-bvb501-isomorphism}{Frame and model isomorphisms}
A frame isomorphism is a bijection \(f:W\to W'\) such that
\(wR_mv\) iff \(f(w)R'_mf(v)\), for every modality and pair of states.
A model isomorphism additionally satisfies
\(w\in V(p)\) iff \(f(w)\in V'(p)\), for every proposition letter.
The second condition cannot be omitted from a truth-preservation claim.
\end{hbdefinition}

\begin{hbtheorem}{thm:p08-bvb501-isomorphism}{proved=annex:p08.src025501.iso}{Source Proposition 2: isomorphism invariance}
If \(f:M\to M'\) is a model isomorphism, then
\(M,w\entails\varphi\) iff \(M',f(w)\entails\varphi\)
for every basic modal formula.
\end{hbtheorem}

\subsection{Evaluation as a finite game}
\label{sec:p08-bvb501-evaluation-games}
\index{games!evaluation}

\begin{hbdefinition}{def:p08-bvb501-evaluation-game}{Verifier and Falsifier}
At a position \((w,\varphi)\), Verifier seeks truth and Falsifier seeks
falsity. At a disjunction Verifier chooses its disjunct; at a conjunction
Falsifier chooses its conjunct. At negation the players exchange roles.
At a diamond Verifier chooses an accessible successor; at a box Falsifier
chooses one. A player required to move without a legal move loses.
At an atom Verifier wins exactly when it is true at the current state;
\(\topv\) and \(\bottom\) have their fixed outcomes. Expand
implication and equivalence by their Boolean definitions for this game.
Every move descends in the formula, so every play terminates.
\end{hbdefinition}

\begin{hbtheorem}{thm:p08-bvb501-evaluation}{proved=annex:p08.src025501.game}{Adequacy of the evaluation game}
Verifier has a winning strategy in the \(\varphi\)-game at \(w\)
exactly when \(M,w\entails\varphi\). Otherwise Falsifier has a winning
strategy. This holds on infinite as well as finite models.
\end{hbtheorem}

\begin{hbexample}{def:p08-bvb501-game-example}{The four-state game in source Figures 5 and 6}
Let \(W=\{1,2,3,4\}\),
\(R=\{(1,3),(1,4),(2,1),(4,2),(4,4)\}\), and
\(V(p)=\{1,2\}\). The complete truth calculation is
\[
\begin{array}{c|cccc}
 &1&2&3&4\\\hline
 \Diamond p&0&1&0&1\\
 \Box\Diamond p&0&0&1&1\\
 \Diamond\Box\Diamond p&1&0&0&1
\end{array}
\]
At state 1, Verifier may choose 3 at the first diamond, leaving Falsifier
no box move. Alternatively she chooses 4; if Falsifier then chooses 4,
she chooses 2 at the final diamond, and if he chooses 2 she chooses 1.
Both terminal states satisfy \(p\). These are the two winning branches
of the source's displayed partial game tree. At state 2 the outer move
must go to 1; Falsifier chooses 3, where the final diamond has no move.
At state 3 the first diamond already has no move. State 4 has the same
winning continuation through 4 as described above.
If Verifier chooses 4 at the final diamond instead of 2, she loses because
\(p\) is false at 4. A winning strategy selects the successful continuation;
it does not make every legal continuation successful.
\end{hbexample}

\section{The first-order standard translation}
\label{sec:p08-bvb501-standard-translation}
\index{standard translation}
\index{correspondence language!first-order}
\index{compactness}
\index{Loewenheim--Skolem property@L\"owenheim--Skolem property}
\index{K}
\index{modus ponens}
\index{modal generalisation}

\begin{hbdefinition}{def:p08-bvb501-standard-translation}{Correspondence language and translation}
Use a binary predicate \(R_m\) for each modality and a unary predicate
\(P_p\) for each proposition letter. Regard the modal model as a
first-order structure with these interpretations. The source translation
\(\operatorname{ST}_x\) commutes with Boolean connectives and constants,
in particular \(\operatorname{ST}_x(\topv)=\topv\), and has clauses
\[
\begin{aligned}
 \operatorname{ST}_x(p)&=P_p(x),&
 \operatorname{ST}_x(\bottom)&=\bottom,\\
 \operatorname{ST}_x(\langle m\rangle\varphi)
 &=\exists y\,(R_m(x,y)\land\operatorname{ST}_y(\varphi)),\\
 \operatorname{ST}_x([m]\varphi)
 &=\forall y\,(R_m(x,y)\to\operatorname{ST}_y(\varphi)).
\end{aligned}
\]
Each recursive quantifier uses a fresh variable. Free variables are contained
in \(\{x\}\); an occurrence of \(x\) need not survive translation
of a Boolean constant. Thus a state assignment supplies the internal modal
perspective without requiring every translated formula to have a free variable.
\end{hbdefinition}

For example,
\(\operatorname{ST}_x(p\to\Diamond p)=
P_p(x)\to\exists y(R(x,y)\land P_p(y))\).
The translation preserves truth at the corresponding assigned state, not
merely the existence of some satisfying model.

\begin{hbtheorem}{thm:p08-bvb501-standard-translation}{proved=annex:p08.src025501.st}{Source Proposition 3: truth under standard translation}
For every model \(M\), state \(w\), and modal formula \(\varphi\),
\[
 M,w\entails\varphi\quad\text{iff}\quad
 M\entails\operatorname{ST}_x(\varphi)[x\leftarrow w].
\]
\end{hbtheorem}

\begin{hbtheorem}{thm:p08-bvb501-modal-metatheory}{proved=annex:p08.src025501.meta}{Source Proposition 4 with its cardinal and effectiveness conditions}
If every finite subset of \(\Sigma\) is locally satisfiable, so is
\(\Sigma\). If \(\Sigma\) has an infinite pointed model, then it has
one of each infinite size \(\kappa\) at least the size of its signature.
For countable signatures this means every infinite cardinality.
For an effectively presented signature the set of basic modal validities is
recursively enumerable.
\end{hbtheorem}

The compactness proof uses one fresh constant for the distinguished state.
The cardinal assertion is about local satisfaction and ordinary cardinal
size, not finite model size or a bound on reachable depth. To see why the
signature qualification matters, for a large set \(J\) use atoms \(p_j\)
and premises \(\Diamond p_j\), together with
\(\Box\neg(p_i\land p_j)\) for every \(i\ne j\).
Any satisfying state has at least \(\abs J\) distinct successors.
This cannot be a model of an arbitrarily smaller infinite cardinality.

For its single-modality case the source also mentions \(K\): propositional tautologies and the
distribution schema, with modus ponens and theorem-only necessitation.
The source's semantic task does not require a premise-necessitation rule.
The registered sheet \ref{sheet:hilbert-k} and the exact canonical proof in
Theorem~\ref{thm:p08-modal-canonical-completeness} supply its theoremhood
axiomatization in the corresponding basic language. These links supply no
proofs for the later richer-language or alternative-semantics claims.

\section{Bisimulation and the distinctions a modal language can draw}
\label{sec:p08-bvb501-bisimulation}
\index{bisimulation}
\index{bisimulation!contraction}
\index{bisimulation!unraveling}
\index{tree unraveling}
\index{unraveling}
\index{bounded morphisms!models}
\index{bounded morphic images!models}
\index{p-morphism}
\index{disjoint unions!models}
\index{generated submodels}
\index{Bisimulation Invariance Lemma}

\begin{hbdefinition}{def:p08-bvb501-bisimulation}{Atomic harmony, zig and zag}
A bisimulation between models \(M\) and \(N\) of the same signature is
a nonempty relation \(E\subseteq W_M\times W_N\). Whenever \(wEw'\),
the two states agree on every proposition letter, and for every modality
\(m\): each \(wR_mv\) has a match \(w'R'_mv'\) with \(vEv'\),
and each \(w'R'_mv'\) has a match \(wR_mv\) with \(vEv'\).
The source names the two transition conditions zig and zag. Pointed models
are bisimilar when such an \(E\) relates their distinguished points.
Unpointed models are bisimilar in the source's existential sense when some
nonempty bisimulation exists; its domain and range need not cover all states.
Modal equivalence means agreement on all finitary basic modal formulas.
\end{hbdefinition}

\begin{hbtheorem}{thm:p08-bvb501-bisimulation-invariance}{proved=annex:p08.src025501.bisim}{Source Lemma 9: bisimulation invariance}
Related states satisfy exactly the same modal formulas. This remains true
when the modal language permits arbitrary set-sized conjunctions and
disjunctions, with formulas built by well-founded recursion.
\end{hbtheorem}

\begin{hbexample}{def:p08-bvb501-three-graphs}{The source's three comparison graphs}
Ignore proposition letters. In the source's Figures~7, 8 and 10, describe the models by
\[
\begin{aligned}
 W_M&=\{s,a,b,d\},&R_M&=\{(s,a),(a,s),(a,b),(s,d)\},\\
 W_N&=\{t,c,b',d'\},&R_N&=\{(t,c),(c,t),(t,b'),(t,d')\},\\
 W_K&=\{u,v\},&R_K&=\{(u,u),(u,v)\}.
\end{aligned}
\]
Both live states of \(M\) have an immediate terminal successor; the
live state \(c\) of \(N\) does not. Hence
\(\Box(\Box\bottom\lor\Diamond\Box\bottom)\) holds at \(s\)
and \(u\), but not at \(t\). The relation
\(E=\{(s,u),(a,u),(b,v),(d,v)\}\) is a bisimulation between
\(M\) and \(K\). Its first designated point is irreflexive and its
second is reflexive. Thus neither \(Rxx\) nor \(\neg Rxx\), as a
property of the current state on arbitrary models, has a basic modal
definition. This is a statement about point-definability, distinct from
frame validity with every valuation quantified over.
\end{hbexample}

\begin{hbexample}{def:p08-bvb501-epistemic-comparison}{Two-state and four-state knowledge diagrams}
The source's Figure~9 compares two universally accessible states, one with
\(p\) and one with \(q\), to four universally accessible states, two
with \(p\) and two with \(q\). Relate precisely states with the same
valuation. Each side has both valuation patterns, so every move has a
matching move. The designated \(p\)-states are bisimilar even though the
models have different sizes. Multiplicity of these alternatives adds no
basic modal information. The same two-state panel appears with the earlier
epistemic illustration.
\end{hbexample}

\begin{hbdefinition}{def:p08-bvb501-transformations}{Unions, generated restrictions and bounded morphisms}
For a nonempty family of models use disjoint tagged copies before taking
their union. Relations and valuations are unions within the components;
there are no transitions between different tags. A generated submodel is
the induced restriction to a nonempty subset closed under every successor
relation. The submodel generated by \(S\ne\varnothing\) contains all
states reachable from \(S\) by finite labelled paths, including paths
of length zero.

A bounded morphism \(f:M\to N\) is a total function satisfying atomic
harmony, relation preservation, and the back condition
\[
 f(w)R'_mv'\quad\Longrightarrow\quad
 \exists v\,(wR_mv\ \text{and}\ f(v)=v').
\]
Its graph is a bisimulation. Surjectivity is the extra condition for
\(N\) to be a bounded morphic image.

Unravel \(M\) at \(w\) into all finite labelled paths from \(w\).
An \(m\)-edge extends a path by one \(R_m\)-step, and its valuation
is that of its last state. Projection to that state is a bounded morphism
onto the submodel generated by \(w\), not necessarily onto an unreachable
part of \(M\).
\end{hbdefinition}

\begin{hbtheorem}{thm:p08-bvb501-transformations}{proved=annex:p08.src025501.transform}{Source Lemma 10 and rooted unraveling}
Modal truth is preserved at corresponding states by disjoint-component
inclusion, generated-submodel inclusion, every bounded morphism, and
rooted unraveling. Arbitrary restrictions need not preserve it.
\end{hbtheorem}

\begin{hbdefinition}{def:p08-bvb501-contraction}{Bisimulation contraction}
Let \(\sim\) be the union of all self-bisimulations of \(M\).
Its quotient classes are written \([w]\); the source notation
\(\lvert w\rvert\) denotes this class here, not its cardinality.
Put \([w]R_m^{\sim}[v]\) iff some representatives are \(R_m\)-related,
and make \(p\) true at \([w]\) iff it is true at \(w\).
\end{hbdefinition}

\begin{hbtheorem}{thm:p08-bvb501-contraction}{proved=annex:p08.src025501.contract}{The contraction is well defined and has no bisimilar duplicate states}
The relation \(\sim\) is an equivalence relation. The quotient valuation
is well defined and its projection is a surjective bounded morphism.
Two distinct quotient states cannot be bisimilar. Consequently no further
surjective bounded morphism can identify two quotient states.
\end{hbtheorem}

For \(M\) above the classes are its two live and two terminal states,
giving exactly \(K\). Unraveling \(K\) at \(u\) produces the comb in
the source's Figure~11: the spine paths are \((u),(u,u),(u,u,u),\ldots\), each with
one terminal path ending in \(v\). The construction expands rather than
contracts the graph, while preserving its root information.

In the source's Figure~12 the outside state \(a\) has edges into the subgraph, but no
path from \(s\) reaches it. Name the bottom and right states \(b,c\)
and the upper terminal state \(d\). The full relation is
\(\{(a,s),(a,b),(b,s),(s,c),(c,b),(s,d)\}\).
Its \(s\)-generated domain is \(\{s,b,c,d\}\).
Deleting \(a\) therefore preserves truth at \(s\).
For contrast, deleting the only successor of a root changes
\(\Diamond\topv\) from true to false. The constant map in the source's Figure~13,
from \((\mathbb N,<)\) onto a single reflexive state, is bounded: an edge
maps to the loop, and each \(n\) can answer the loop with \(n+1\).

\subsection{First-order structural comparison}
\label{sec:p08-bvb501-first-order-comparison}

The source compares its modal structural equivalence with the choice of
invariants in Klein's geometric program. Different questions require
different equivalences. For first-order formulas an isomorphism preserves
truth of every finite tuple, but agreement on sentences does not generally
determine cardinality or isomorphism type.

\begin{hbdefinition}{def:p08-bvb501-potential-isomorphism}{Potential isomorphism and infinitary formulas}
A finite partial isomorphism \(f:M\to N\) is an injective finite map
preserving and reflecting every atomic formula on tuples from its domain,
including atoms with constants or terms in the signature. A potential
isomorphism is a nonempty family \(I\) of such maps with forth and back:
for each \(f\in I\) and \(a\in M\), some extension in \(I\)
includes \(a\) in its domain, and symmetrically for each \(b\in N\)
in the range. Replacing \(I\) by its restriction closure preserves these
conditions.

The language \(L_{\infty\omega}\) permits arbitrary set-sized
conjunctions and disjunctions with only finitely many free variables and
finite quantifier strings. Formulas are well founded. Its countable-connective
fragment is \(L_{\omega_1\omega}\). These are first-order infinitary
languages; the infinitary modal language below instead extends modal
connectives and has no world variables in its syntax.
\end{hbdefinition}

\begin{hbtheorem}{thm:p08-bvb501-potential-isomorphism}{proved=annex:p08.src025501.potential}{Potential isomorphism, infinitary theories and countable definitions}
Potentially isomorphic structures agree on every
\(L_{\infty\omega}\)-formula at tuples matched by a map in the family.
Every structure \(M\) has a sentence \(\delta_M\) of that language
whose models are exactly the structures potentially isomorphic to \(M\).
Consequently equality of complete \(L_{\infty\omega}\)-theories is
equivalent to potential isomorphism.
If \(M\) and its signature are countable, \(\delta_M\) can lie in
\(L_{\omega_1\omega}\) and characterizes \(M\) up to isomorphism
among countable structures.
\end{hbtheorem}

\begin{hblemma}{thm:p08-bvb501-first-order-game}{proved=annex:p08.src025501.fo-game}{Finite-rank first-order comparison}
For a finite relational signature, allowing finitely many constants but
no function symbols, Duplicator wins the \(k\)-round first-order
comparison game at matching finite tuples iff those tuples agree on every
formula of quantifier rank at most \(k\). A round chooses an arbitrary
domain element in either structure and answers in the other; atomic
disagreement loses immediately. The finite-signature condition belongs
to this finitary strategy-to-formula converse.
\end{hblemma}

For \((\mathbb Q,<)\) and \((\mathbb R,<)\), all finite increasing
tuple matches form a potential isomorphism: a new point is placed in the
same interval or ray on the other side, where density and absence of
endpoints supply a fresh point. These structures have the same first-order
theory, but are not isomorphic because their cardinalities differ.
For finite structures elementary equivalence does imply isomorphism:
finitely many candidate bijections can be excluded by a finite atomic diagram
if none is an isomorphism. More generally, a logic with a set of formulas has
only a set of complete theories, while structures of arbitrarily different
cardinalities form a proper class of isomorphism types. Isomorphism invariance
alone cannot give the converse for every structure.

\begin{hbdefinition}{def:p08-bvb501-assignment-model}{Turning assignments into modal states}
For a first-order structure \(M\), take all assignments \(g\) of its infinite
variable set to its domain as states. For each variable \(x\), the
\(x\)-transition resets only that variable to an arbitrary domain element.
Give each atomic first-order formula its own proposition letter, true at
an assignment exactly when that atomic formula is true there.
This modal signature can be infinite, even for a finite first-order signature.
\end{hbdefinition}

\begin{hbtheorem}{thm:p08-bvb501-assignment-bridge}{proved=annex:p08.src025501.assignments}{Potential isomorphism as existential assignment-model bisimilarity}
Two first-order structures are potentially isomorphic iff their assignment
models have a nonempty bisimulation. This existential statement does not
say every assignment on either side can be matched; finite-support
assignments suffice for the constructed witness.
\end{hbtheorem}

\subsection{Modal equivalence and its converse hypotheses}
\label{sec:p08-bvb501-equivalence}

\begin{hbtheorem}{thm:p08-bvb501-hennessy-milner}{proved=annex:p08.src025501.hm}{Source Proposition 11 and its image-finite extension}
If each state of \(M\) and \(N\) has finitely many successors for
each modality, modal equivalence is itself a bisimulation wherever it holds.
This includes finite models and permits infinite proposition signatures.
\end{hbtheorem}

The finite successor condition cannot be removed. The source's Figure~14 has a root
with one finite terminal chain of each length, compared to the same model
with an additional infinite chain from the root. All valuations are empty.
Every formula has finite modal depth and inspects only a bounded initial
part of a chosen chain. Both roots satisfy the same formulas. A bisimulation
cannot match the infinite chain to any finite one, because successive
matching steps would eventually face a terminal state on only one side.
The comparison-game proof below makes both assertions explicit.

\begin{hbdefinition}{def:p08-bvb501-modal-saturation}{Realizing successor types}
A model is modally saturated when, for every state \(w\), modality \(m\)
and set \(\Gamma\) of modal formulas, if every finite subset of
\(\Gamma\) holds at some \(m\)-successor of \(w\), all of
\(\Gamma\) holds at one such successor. Successor sets, not arbitrary
global realizations, enter this condition. An \(\omega\)-saturated
first-order structure realizes every complete finite-tuple type over every
finite parameter set consistent with its elementary diagram.
\end{hbdefinition}

\begin{hbtheorem}{thm:p08-bvb501-first-lifting}{proved=annex:p08.src025501.lift1}{Source Lemma 15: elementary extensions become bisimilar}
Every relational first-order structure has an \(\omega\)-saturated
elementary extension. Such an extension of a modal correspondence structure
is modally saturated. Two pointed modal models are modally equivalent iff
they have pointed elementary extensions that are bisimilar.
\end{hbtheorem}

\begin{hbtheorem}{thm:p08-bvb501-characterization}{proved=annex:p08.src025501.character}{Source Theorem 13: the modal fragment of first-order logic}
A first-order correspondence formula \(\alpha(x)\) is equivalent on
all models to \(\operatorname{ST}_x(\varphi)\) for a basic modal
formula iff its truth is invariant under pointed bisimulations.
This is the all-model theorem. Restricting to finite models removes the
compactness argument used in its proof and requires a separate proof.
\end{hbtheorem}

\subsection{Interpolation and the two lifting routes}
\label{sec:p08-bvb501-interpolation}
\index{Craig interpolation}
\index{Modal Characterisation Theorem}

\begin{hbdefinition}{def:p08-bvb501-relative-entailment}{Entailment along a common-signature bisimulation}
Let \(\alpha(x)\) and \(\beta(x)\) be first-order correspondence
formulas, and let \(L\) be their common predicate and accessibility
signature. Say that \(\alpha\) implies \(\beta\) along
\(L\)-bisimulations when, whenever pointed models of their respective
full signatures have bisimilar \(L\)-reducts, truth of \(\alpha\)
at the first point implies truth of \(\beta\) at the second.
Only atoms and modalities in \(L\) are checked by this bisimulation.
The restriction is essential: a relation checking private symbols would
not express the common-language interpolation requirement.
\end{hbdefinition}

\begin{hbtheorem}{thm:p08-bvb501-relative-interpolation}{proved=annex:p08.src025501.interpolate}{Source Theorem 14 with the common-language condition explicit}
\(\alpha\) implies \(\beta\) along \(L\)-bisimulations iff
there is a basic modal \(L\)-formula \(\theta\) such that
\(\alpha\entails\operatorname{ST}_x(\theta)\) and
\(\operatorname{ST}_x(\theta)\entails\beta\).
All models are allowed. In the equal-signature case \(\alpha=\beta\),
this recovers the modal characterization theorem.
\end{hbtheorem}

The source refers the interpolation proof to Barwise and van Benthem,
\emph{Interpolation, Preservation, and Pebble Games},
\emph{Journal of Symbolic Logic} 64(2) (1999), 881--903.
The common-reduct reading also appears explicitly in van Benthem's
advanced-modal-logic course formulation \cite{VanBenthem2005AMLNotes}. Entailment along a structural
relation compares different models, whereas ordinary entailment starts
with a single model. Thus ordinary modal interpolation needs an additional
argument; it does not follow just by replacing \(\alpha,\beta\)
with standard translations in the preceding statement.

\begin{hbtheorem}{thm:p08-bvb501-modal-interpolation}{proved=annex:p08.src025501.modal-interpolate}{Ordinary modal interpolation over a fixed modality signature}
Fix the accessibility modalities. If basic modal formulas
\(\varphi\entails\psi\), there is a basic modal formula
\(\theta\), using only proposition letters common to
\(\varphi,\psi\), with \(\varphi\entails\theta\entails\psi\).
The proof uses a model on related pairs to combine private valuations.
The modalities remain fixed; they are not being treated as private
proposition letters.
\end{hbtheorem}

\begin{hbtheorem}{thm:p08-bvb501-second-lifting}{cited=BlackburnVanBenthem2006ILLC}{Source Lemma 16: bisimilar covers become elementarily equivalent}
Two pointed modal models are modally equivalent iff each is bisimilar to
a pointed model such that the resulting two pointed correspondence
structures are elementarily equivalent. Elementary equivalence here
uses a named constant for the distinguished point.
\end{hbtheorem}

The first lifting theorem changes the models by elementary extensions
and produces a bisimulation. The second changes them within their
bisimulation classes and produces elementary equivalence. Its source
proof route is tree unraveling, duplication of subtrees, and a
first-order back-and-forth argument; the manuscript refers the actual
construction to van Benthem's \emph{Bisimulation: the Never Ending Story}
(1996), pp.~23--27, and Andr\'eka, van Benthem and N\'emeti's
\emph{Modal Logics and Bounded Fragments of Predicate Logic}
(1998), pp.~217--274. Plain unraveling alone is not the asserted
elementary-equivalence construction, and the first lifting proof does
not prove the second result. In the reverse implication, elementary
equivalence preserves every standard translation, and the two
bisimulations preserve every modal formula.

\begin{hbtheorem}{thm:p08-bvb501-finite-characterization}{cited=BlackburnVanBenthem2006ILLC}{The finite-model characterization reported in source Section 3.5}
For a fixed finite correspondence vocabulary, a first-order formula
invariant under bisimulations between finite pointed models is equivalent
on finite pointed models to a basic modal formula. The source reports
Rosen's finite-model theorem and Otto's refinement giving a modal-depth
bound \(2^k\) when the first-order quantifier depth is \(k\).
These are finite-model results, not applications of first-order compactness.
\end{hbtheorem}

Rosen's \emph{Modal Logic over Finite Structures} (1997), pp.~427--439,
and Otto's 2004 Darmstadt report \emph{Elementary Proof of the van
Benthem--Rosen Characterisation Theorem} are the proof references identified
by the source. The source supplies no finite-model proof here.
Its broader analogy between modal and first-order metaproperties is
a research perspective: compactness, interpolation and cardinal-transfer
phenomena have modal counterparts, but the analogy is not a theorem
that every first-order metaproperty transfers, and decidability already
separates the two languages.

\subsection{Comparison games and their precise finitary conditions}
\label{sec:p08-bvb501-comparison-games}
\index{games!bisimulation}

\begin{hbdefinition}{def:p08-bvb501-comparison-game}{Spoiler and Duplicator}
Start at a pair of states. Atomic disagreement is an immediate Spoiler win.
In a round Spoiler chooses a modality and a successor in either model;
Duplicator must choose a successor for the same modality in the other.
An impossible response or atomic disagreement is a Spoiler win. If Spoiler
has no legal move Duplicator wins. In the \(k\)-round game Duplicator
wins after surviving \(k\) rounds; in the unlimited game she wins an
infinite play without a loss as well. Restricting to a relevant signature
restricts both the atomic checks and the available modality choices.
\end{hbdefinition}

\begin{hbtheorem}{thm:p08-bvb501-infinite-game}{proved=annex:p08.src025501.game-infinite}{Unlimited games, bisimulations and determinacy}
Duplicator wins the unlimited comparison game iff the starting points are
bisimilar. Otherwise Spoiler has a winning strategy. The statement holds
for arbitrary set-sized models and signatures; winning plays may have no
uniform finite length bound.
\end{hbtheorem}

\begin{hbtheorem}{thm:p08-bvb501-finite-game}{proved=annex:p08.src025501.game-finite}{The valid finitary forms of source Lemma 17}
For a finite relevant proposition and modality signature, Spoiler wins the
\(k\)-round game iff a formula of modal depth at most \(k\)
distinguishes the starting points. The same equivalence holds for arbitrary
signatures when both models are image-finite for every modality. Under
either condition Duplicator wins iff all formulas of depth at most \(k\)
agree. Without these conditions the implication from a distinguishing
formula to a Spoiler strategy still holds; its converse can fail.
\end{hbtheorem}

\begin{hbexample}{def:p08-bvb501-game-counterexample}{A one-round win without a finitary distinguishing formula}
Use one relation and atoms \(p_i\), \(i<\omega\). Model \(M\)
has a root with one terminal successor for every finite subset
\(F\subseteq\omega\), making precisely those \(p_i\) with
\(i\in F\) true. Model \(N\) has all those successors and one more
terminal successor making every \(p_i\) true. Both roots make all atoms
false. Any finitary formula uses a finite atom set \(A\); both successor
families realize every \(A\)-pattern. At terminals, modal truth depends
only on that pattern and the vacuous modal clauses. Induction then gives
identical root truth for every finitary formula.

Spoiler chooses the all-atoms successor in \(N\) on the first round.
Every response in \(M\) has a finite \(F\) and disagrees on some
\(p_i\), \(i\notin F\). Thus Spoiler wins in one round although
there is no finitary distinguishing formula of any modal depth. This
exhibits the missing conditions directly; finite-round strategy adequacy
is not inferred from the unlimited-game equivalence.
\end{hbexample}

For the earlier \(M,N,K\) graphs, the displayed depth-three formula gives
Spoiler a win within three rounds. Duplicator wins between \(M,s\) and
\(K,u\) by maintaining the explicit relation \(E\). For the finite-chain
and infinite-chain example, she survives every fixed finite number of rounds
by initially choosing a sufficiently long finite chain, yet loses the unlimited
game when Spoiler commits to the infinite chain. These are distinct games.

\begin{hbdefinition}{def:p08-bvb501-infinitary-modal-language}{The source's infinitary modal notation}
Extend the modal grammar by arbitrary set-sized conjunctions and disjunctions,
using well-founded formulas. The source also calls this language
\(L_{\infty\omega}\), now in a modal role rather than the first-order
role of Definition~\ref{def:p08-bvb501-potential-isomorphism}.
Its formulas have no world variables or first-order quantifiers.
Restricting the connective indices to countable sets gives the corresponding
countable-infinitary modal language.
\end{hbdefinition}

\begin{hbtheorem}{thm:p08-bvb501-infinitary-modal}{proved=annex:p08.src025501.modal-infinitary}{Infinitary modal definition up to bisimulation}
For every pointed model there is a set-sized infinitary modal formula
characterizing it up to pointed bisimulation. For a countable model and
countable signature, the formula can use only countable conjunctions and
disjunctions. Consequently, two such pointed countable models have the same
infinitary modal theory iff they are bisimilar.
\end{hbtheorem}

\section{Computation and reasoning tasks}
\label{sec:p08-bvb501-computation}
\index{model checking}
\index{labelling algorithm}
\index{global satisfiability problem}
\index{TBox}
\index{ABox}

\subsection{Finite-model checking and representation}
\label{sec:p08-bvb501-model-checking}

\begin{hbmethod}{def:p08-bvb501-labelling}{Bottom-up truth-set computation}
For an explicitly represented finite model, compute the truth set of every
subformula once, in subformula order. Atoms use their valuations; Boolean
connectives use complement, intersection, union and the corresponding
Boolean set operations. For a modality use
\[
 \begin{aligned}
 [\![\langle m\rangle\sigma]\!]&=
  \{w:\exists v\,(wR_mv\land v\in[\![\sigma]\!])\},\\
 [\![[m]\sigma]\!]&=
  \{w:\forall v\,(wR_mv\Rightarrow v\in[\![\sigma]\!])\}.
 \end{aligned}
\]
The brackets here denote the computed set of satisfying states. The
answer to local checking is membership of the designated point; global
checking asks whether the answer set is the entire domain.
\end{hbmethod}

\begin{hbtheorem}{thm:p08-bvb501-model-checking}{proved=annex:p08.src025501.check}{Correctness and explicit-graph running time}
The bottom-up algorithm computes exactly the satisfying states of a basic
modal formula. With adjacency lists and valuation bit sets its time is
\(O(s(\varphi)(|W|+\sum_m|R_m|))\), counting only modalities
occurring in the formula; \(s(\varphi)\) is its number of subformulas.
For one relation represented by a matrix this gives the source's
\(O(s(\varphi)|W|^2)\) bound. Store space is
\(O(s(\varphi)|W|)\), besides the input.
\end{hbtheorem}

A model lists facts and links, while a nested formula asks whether a
structured condition follows from those facts. This makes model checking
a reasoning task, distinct from constructing a model or proving validity.
The source's hardware example uses a graph of possible runs and a modal
specification of permitted behavior. Which run properties are expressible
depends on the chosen language; basic one-step modalities alone do not
express arbitrary temporal requirements.

\begin{hbdefinition}{def:p08-bvb501-until-check}{Path Until and interval Until}
On a directed graph define \(U_{\exists}(\psi,\theta)\) at \(s\)
to mean that a finite path \(s=w_0R\cdots Rw_n\) ends at a
\(\psi\)-state and every \(w_i\), \(i<n\), satisfies
\(\theta\). Length zero is allowed. The source's Figure~15 algorithm
initially marks the \(\psi\)-states, puts them in a queue, and repeatedly
marks an unmarked predecessor satisfying \(\theta\), adding it to the
queue. Each state enters the queue at most once.

The source also prints an interval clause: there is a \(t\) with
\(sR^*t\) satisfying \(\psi\), and every \(u\) with
\(sR^*u\) and \(uR^+t\) satisfies \(\theta\).
Here \(R^*\) and \(R^+\) are reflexive-transitive and transitive
closure. On a graph with branching and merging, that universal interval
clause is stronger than existence of a suitable path. On a directed
acyclic graph having a unique path between any related pair, the interval
is exactly the chosen path before its endpoint, so the clauses coincide.
\end{hbdefinition}

\begin{hbtheorem}{thm:p08-bvb501-until-check}{proved=annex:p08.src025501.until}{What the source queue algorithm computes}
The queue algorithm computes exactly \(U_{\exists}(\psi,\theta)\)
in \(O(|W|+|R|)\) time once the two argument truth sets are known.
Consequently it also computes the printed interval clause on directed
acyclic unique-path graphs, including successor chains. It does not
compute that interval clause on arbitrary irreflexive graphs.
\end{hbtheorem}

\begin{hbexample}{def:p08-bvb501-until-counterexample}{A merging branch separates the two clauses}
Take edges \(sRa,aRt,sRb,bRt\), make \(\psi\) true only
at \(t\), and make \(\theta\) true at \(s,a\) but false
at \(b\). The graph is finite and irreflexive. The algorithm marks
\(t,a,s\), witnessing the path \(s,a,t\). The interval clause
fails at \(s\), since \(sR^*b\) and \(bR^+t\) but
\(b\) fails \(\theta\). Cycles give an additional reason that
irreflexivity alone does not turn the interval into a finite linear path.
\end{hbexample}

Explicit graphs and symbolic graphs have different input sizes. The source
discusses binary decision diagrams, their ordered form, and symbolic model
checking as methods for representing Boolean state sets and transition
relations compactly. A graph-polynomial bound does not itself imply a
polynomial bound in a succinct symbolic description. The source's
historical reference chain is Lee and Akers for BDDs, Bryant for ordered
BDDs, and Burch, Clarke, McMillan, Dill and Hwang for symbolic checking.
The number of semantic states a representation describes is not its
stored byte count.

\subsection{Two-variable translation and finite models}
\label{sec:p08-bvb501-filtration}
\index{two-variable fragment}
\index{filtration}
\index{filtration!minimum}
\index{filtration!maximum}
\index{Filtration Theorem}
\index{finite model property}
\index{finite model property!strong}
\index{decidability}

\begin{hbtheorem}{thm:p08-bvb501-two-variable}{proved=annex:p08.src025501.two-variable}{The alternating-variable translation}
Every basic modal formula has a truth-preserving first-order translation
using only variables \(x,y\). Its modal clauses are
\[
 \begin{aligned}
 \operatorname{ST}_x(\langle m\rangle\sigma)&=
   \exists y(R_mxy\land\operatorname{ST}_y(\sigma)),\\
 \operatorname{ST}_y(\langle m\rangle\sigma)&=
   \exists x(R_myx\land\operatorname{ST}_x(\sigma)),\\
 \operatorname{ST}_x([m]\sigma)&=
   \forall y(R_mxy\to\operatorname{ST}_y(\sigma)),\\
 \operatorname{ST}_y([m]\sigma)&=
   \forall x(R_myx\to\operatorname{ST}_x(\sigma)).
 \end{aligned}
\]
The Boolean and atomic clauses are unchanged. Thus modal satisfiability
reduces to satisfiability of a two-variable sentence by existentially
closing the displayed free variable. The source cites two-variable
first-order decidability for this route; the finite-model route below
provides a modal decision proof directly.
\end{hbtheorem}

\begin{hbdefinition}{def:p08-bvb501-paired-filtration}{Filtration for both primitive modalities}
Let \(\Sigma\) be subformula closed. Put \(w\sim_\Sigma v\)
iff \(w,v\) agree on every member of \(\Sigma\); write
\([w]_\Sigma\), or the source's \(|w|_\Sigma\), for a class.
The quotient has these classes as states and the inherited valuation for
atoms in \(\Sigma\). A filtration relation \(R_m^f\) must obey:
every original \(R_m\)-edge induces a quotient edge; and every quotient
edge from \([w]_\Sigma\) to \([v]_\Sigma\) satisfies both
\[
 \begin{aligned}
 \langle m\rangle\sigma\in\Sigma, M,v\entails\sigma
     &\Rightarrow M,w\entails\langle m\rangle\sigma,\\
 [m]\sigma\in\Sigma, M,w\entails[m]\sigma
     &\Rightarrow M,v\entails\sigma.
 \end{aligned}
\]
All these conditions depend only on the classes. The minimum relation
\(R_m^s\) contains edges induced by some original edge; the maximum
\(R_m^l\) contains exactly class pairs satisfying both displayed
conditions. Any relation between these two is a filtration relation.
The box condition is needed because box is primitive in the source grammar.
\end{hbdefinition}

\begin{hbtheorem}{thm:p08-bvb501-filtration}{proved=annex:p08.src025501.filtration}{Source Theorem 20 with both modal clauses supplied}
Every paired filtration preserves truth of every \(\sigma\in\Sigma\)
at its class. The minimum and maximum relations are well defined,
\(R_m^s\subseteq R_m^l\), and every intervening relation works.
For finite \(\Sigma\) the quotient has at most \(2^{|\Sigma|}\)
states.
\end{hbtheorem}

\begin{hbexample}{def:p08-bvb501-filtration-counterexample}{Why a diamond-only constraint is insufficient}
Let \(M\) have two isolated states \(w,v\), with \(p\) true
only at \(w\), and let \(\Sigma=\{p,\Box p\}\).
Both states satisfy \(\Box p\), and their classes are distinct.
The printed diamond-only maximum imposes no constraint, because
\(\Sigma\) contains no diamond, so it allows an edge
\([w]_\Sigma R^f[v]_\Sigma\). Then \(\Box p\) becomes
false at \([w]_\Sigma\). The paired condition excludes precisely
this edge. The original minimum filtration already works; the erroneous
unrestricted maximum does not preserve the full primitive language.
\end{hbexample}

\begin{hbcorollary}{thm:p08-bvb501-finite-decision}{proved=annex:p08.src025501.finite-decision}{Source Theorem 18 and a terminating decision method}
If \(\varphi\) is satisfiable, it has a model with at most
\(2^{s(\varphi)}\) states. Satisfiability and validity of the basic
modal language are decidable: enumerate all models up to that bound
over the finite relevant atom and modality signatures and check truth.
No bound on the ambient signature is needed, since symbols not occurring
in the input formula can be ignored or interpreted arbitrarily.
\end{hbcorollary}

First-order logic can force infinity: the conjunction asserting
irreflexivity, transitivity and \(\forall x\exists y\,Rxy\)
does so. Starting at any point and following serial edges gives an
arbitrarily long chain of distinct points: a repeated point would create
a cycle, and transitivity would contradict irreflexivity. The positive
integers with strict order satisfy these conditions. Basic modal
satisfiability has no corresponding infinity requirement because of its
finite-model bound. Filtration is one decision method; selective witness
construction, tableaux and automata provide other routes, with different
representations and bounds.

\subsection{Depth, space and complexity boundaries}
\label{sec:p08-bvb501-complexity}

\begin{hbtheorem}{thm:p08-bvb501-tree-depth}{proved=annex:p08.src025501.tree-depth}{Source Lemma 21 and the satisfiable form of Theorem 22}
Unravel a pointed model into a rooted tree and truncate it after depth
\(d\). The root's formulas of modal depth at most \(d\) keep their
truth values. Every satisfiable basic modal formula \(\varphi\)
therefore has a finite rooted tree model of depth at most
\(\operatorname{md}(\varphi)\). The word ``satisfiable'' is essential;
an inconsistent formula has no model of any depth.
\end{hbtheorem}

The finite-model bound permits nondeterministically guessing an
exponential-size model and checking it in exponential time. Depth-first
search through modal obligations uses less space: it need not keep all
branches of a satisfying tree simultaneously.

\begin{hbtheorem}{thm:p08-bvb501-space-search}{proved=annex:p08.src025501.space-search}{A polynomial-space modal search}
Basic modal satisfiability has a nondeterministic decision algorithm using
polynomial space in the formula length. It stores a stack of Boolean-
consistent local obligations; each diamond demands one recursively
checked successor and each box constrains every such successor. Stack
depth is bounded by modal depth. With the standard equality
\(\mathrm{NPSPACE}=\mathrm{PSPACE}\), this is a deterministic
polynomial-space upper bound.
\end{hbtheorem}

\begin{hbtheorem}{thm:p08-bvb501-complexity-classifications}{cited=BlackburnVanBenthem2006ILLC}{Complexity classifications reported by the source}
Basic modal satisfiability is PSPACE-complete. Full first-order model
checking has PSPACE-complete combined complexity, while the basic modal
explicit-model checking algorithm is polynomial. Two-variable first-order
satisfiability is NEXPTIME-complete. These are distinct input problems;
the two-variable reduction supplies an upper bound but does not identify
the modal problem with the full two-variable problem.
\end{hbtheorem}

The source explicitly defers a detailed PSPACE algorithm and hardness
proof. The depth-first upper-bound construction is supplied here; the
hardness and the other classifications remain the source's cited results.
PSPACE-completeness does not prove \(\mathrm P\ne\mathrm{PSPACE}\).
It supports a conditional warning about worst-case polynomial time,
not an unconditional claim that every useful modal instance is difficult.

\subsection{Global constraints and comparison}
\label{sec:p08-bvb501-other-tasks}

Global satisfiability asks whether \(\psi\) is true at every point of
some nonempty model. Local--global satisfiability asks for such a model
with a selected point satisfying \(\varphi\). A background theory
can impose a grammar of permitted structures while a local query asks
whether a particular configuration can occur. This distinction is used
in description-logic TBoxes and queries.

\begin{hbexample}{def:p08-bvb501-bear-care}{The source's BEAR-CARE background theory}
Require everywhere
\[
 \begin{gathered}
 \mathrm{bear}\lor\mathrm{human},\qquad
 \mathrm{bear}\to\neg\mathrm{human},\\
 \mathrm{bear}\to\langle\mathrm{MOTHER}\rangle\mathrm{bear},\qquad
 \mathrm{bear}\to[\mathrm{FEDBY}](\mathrm{zoo\mbox{-}keeper}\lor\mathrm{mother}).
 \end{gathered}
\]
No point can satisfy
\(\langle\mathrm{MOTHER}\rangle(\mathrm{bear}\land\mathrm{human})\),
because no successor can be both a bear and human. The query
\(\langle\mathrm{FEDBY}\rangle(\neg\mathrm{human}\land\neg\mathrm{mother})\)
is satisfiable: take two bear states \(s,t\), neither human nor mother,
make \(t\) a zoo-keeper, give each a MOTHER self-loop, and give
\(s\) the sole FEDBY edge to \(t\). Both satisfy every global
axiom and \(s\) satisfies the query. The background theory did not
say that zoo-keepers must be human.
\end{hbexample}

\begin{hbtheorem}{thm:p08-bvb501-global-complexity}{cited=BlackburnVanBenthem2006ILLC}{Source Theorem 23: the cost of global restrictions}
Global and local--global satisfiability for the basic modal language are
EXPTIME-complete. The source refers the proof to Hemaspaandra's
\emph{The Price of Universality} and description-logic literature.
\end{hbtheorem}

This worst-case classification differs from local satisfiability. It does
not give an exact exponential lower bound for every instance or every
algorithm. Practical success on restricted or structured inputs is
compatible with such a classification. The source's finite-graph
comparison task, in contrast, has a direct polynomial algorithm.

\begin{hbtheorem}{thm:p08-bvb501-finite-bisimulation-algorithm}{proved=annex:p08.src025501.finite-bisim-algorithm}{Polynomial finite-model bisimulation comparison}
For two explicitly represented finite models with a finite relevant
signature, start with the atomically agreeing pairs. Repeatedly discard
any pair for which some successor on one side has no remaining matching
successor on the other, for the same modality. The stabilized relation
is the greatest bisimulation, so membership of the designated pair decides
pointed bisimilarity. There are at most \(|W_M||W_N|\) deletion rounds;
a direct implementation is polynomial in the explicit input size.
\end{hbtheorem}

The source compares this with graph isomorphism to illustrate how changing
the required structural relation changes a reasoning task. Its statements
about unresolved graph-isomorphism complexity are historical statements
from 2006, not an undated current classification. Likewise, the source's
selected application examples motivate the algorithms but do not erase
the distinction between expressivity, model checking, satisfiability,
global constraints and model comparison.

\section{Richer logics and frame correspondence}
\label{sec:p08-bvb501-frame-logics}
\index{frame}
\index{frame validity}
\index{validity!at a frame point}
\index{validity!on a frame}
\index{validity!on a class of frames}
\index{normal modal logics}
\index{frame incompleteness}
\index{K4}
\index{S5}

This section fixes one accessibility relation and its box and diamond.
Its normal-logics and partial-function claims use that single-modality
standing; they are not assertions about every multimodal combination.

\subsection{Axiom extensions and their semantic limits}
\label{sec:p08-bvb501-axioms}

\begin{hbdefinition}{def:p08-bvb501-frame-validity}{Validity quantifies over valuations}
For a frame \(F\), write \(F,w\entails\varphi\) when
\((F,V),w\entails\varphi\) for every valuation \(V\), and
\(F\entails\varphi\) when this holds at every point. A formula,
or a set of formulas, defines the frames validating it. A normal modal
logic is a set of formulas containing K and closed under uniform
substitution, modus ponens and theorem necessitation. Inclusion orders
these logics; intersection is their meet and normal closure of union
their join. Syntactic closure alone does not assert completeness for
some class of relational frames.
\end{hbdefinition}

K4 adds all instances of \(\Box\varphi\to\Box\Box\varphi\);
KAlt1 adds all instances of \(\Diamond\varphi\to\Box\varphi\).
Their schemas express transitivity and partial functionality, respectively.
The schemas are not valid on all frames, although special instances with
tautological arguments can already be K-theorems. Frame conditions may
motivate an axiomatization, or a given proof system may motivate a search
for its characteristic frames. GL, for example, adds the L\"ob schema
and has a provability interpretation as well as a relational one.

\begin{hbtheorem}{thm:p08-bvb501-k4-completeness}{cited=BlackburnVanBenthem2006ILLC}{Source Theorem 24}
K4 proves exactly the formulas valid on all transitive frames. GL proves
exactly the formulas valid on finite transitive irreflexive rooted trees,
where a tree's accessibility relation is its strict ancestor relation,
not merely its immediate-child relation.
\end{hbtheorem}

The source cites canonical models and model transformations as different
completeness methods; it does not give these completeness proofs here.
It also reports uncountably many normal modal logics and Makinson's
1971 embedding theorem:

\begin{hbtheorem}{thm:p08-bvb501-makinson}{cited=BlackburnVanBenthem2006ILLC}{Source Theorem 25: the two maximal consistent extensions}
Every consistent normal modal logic is contained in either Id, generated
by \(\varphi\leftrightarrow\Box\varphi\), or Un, generated
by \(\Box\bottom\). Id is the logic of isolated reflexive points,
and Un the logic of frames with no edges. The consistency restriction
excludes the normal logic containing every formula.
\end{hbtheorem}

\begin{hbdefinition}{def:p08-bvb501-tmeq}{The source's incompleteness example, with the independent argument retained}
TMEQ is the normal extension generated by
\[
 \begin{gathered}
 \varphi\to\Diamond\varphi,\qquad
 \Box\Diamond\varphi\to\Diamond\Box\varphi,\\
 \Diamond(\Diamond\varphi\land\Box\psi)
    \to\Box(\Diamond\varphi\lor\Box\psi),\qquad
 (\Diamond\varphi\land\Box(\varphi\to\Box\varphi))\to\varphi.
 \end{gathered}
\]
The last box in the E schema has argument \(\psi\), not
\(\varphi\). This independent argument is also printed in the
TMEQ recollection and appendix of Ding and Li's 2024 author manuscript
\cite{DingLi2024PQMLv1}. The T schema is equivalent, under uniform
substitution and classical duality, to \(\Box\varphi\to\varphi\).
\end{hbdefinition}

\begin{hbtheorem}{thm:p08-bvb501-tmeq-incomplete}{cited=DingLi2024PQMLv1}{The relational-frame incompleteness example}
The frames validating TMEQ have identity accessibility, but TMEQ does
not prove every formula valid on those frames. In particular it does
not prove \(p\to\Box p\). Consequently no class of relational
frames has exactly TMEQ as its logic. This is the incompleteness
phenomenon identified by the source's reference to van Benthem (1978).
\end{hbtheorem}

\begin{hbtheorem}{thm:p08-bvb501-printed-e-collapse}{proved=annex:p08.src025501.e-collapse}{Why identifying the two E arguments changes the theorem}
Replacing the final \(\psi\) in E by \(\varphi\), while
retaining the other schemas, generates Id instead. That extension is
complete for identity frames, so it is not the TMEQ incompleteness example.
\end{hbtheorem}

\subsection{Defining and failing to define frames}
\label{sec:p08-bvb501-frame-definability}
\index{frame definability}
\index{defining a frame class}
\index{disjoint unions!frames}
\index{generated subframes}
\index{bounded morphisms!frames}
\index{bounded morphic images!frames}
\index{filter}
\index{proper filter}
\index{ultrafilter}
\index{ultrafilter!principal}
\index{ultrafilter extensions}

\begin{hbtheorem}{thm:p08-bvb501-elementary-correspondents}{proved=annex:p08.src025501.correspondents}{All four source Proposition 27 correspondents, and reflexivity}
The respective frame correspondents are
\[
 \begin{array}{ll}
 \Box p\to\Box\Box p & R\text{ transitive},\\
 \Diamond p\to\Box p & \forall xyz(Rxy\land Rxz\to y=z),\\
 p\leftrightarrow\Box p & \forall xy(Rxy\leftrightarrow x=y),\\
 \Box\bottom & \forall xy\neg Rxy,\\
 p\to\Diamond p & \forall xRxx.
 \end{array}
\]
Identity accessibility is stronger than reflexivity plus antisymmetry;
the latter still allows distinct related points.
\end{hbtheorem}

\begin{hbtheorem}{thm:p08-bvb501-frame-transformations}{proved=annex:p08.src025501.frame-transform}{Source Theorem 28: three validity-preserving transformations}
Basic modal frame validity is preserved by nonempty disjoint unions,
nonempty generated subframes, and surjective bounded morphic images.
Any class defined by basic modal formulas must therefore be closed
under these operations.
\end{hbtheorem}

These yield the source's three undefinability examples. The class
\(\forall xy(Rxy\lor Ryx)\) is not closed under unions of
separate nonempty components. Frames possessing an isolated reflexive
point are not closed under generated subframes: select the isolated
irreflexive point in a two-component example. Irreflexive frames are
not closed under bounded morphic images: collapse \((\mathbb N,<)\)
onto a singleton loop. None of these classes is basic modally definable,
even by a set of formulas. Failure to define pointwise irreflexivity
in models and failure to define irreflexive frames are different claims;
the latter uses frame-validity preservation over all valuations.

\begin{hbdefinition}{def:p08-bvb501-ultrafilter-extension}{Source Definition 29}
For a frame \(F=(W,R)\), put
\(l(X)=\{w:\forall v(wRv\Rightarrow v\in X)\}\).
A proper filter on \(W\) contains \(W\), is upward closed and
closed under finite intersections, and excludes the empty set.
An ultrafilter decides each subset against its complement. Write
\(\operatorname{uf}(W)\) for these ultrafilters. The extension
\(\operatorname{ue}(F)\) has this domain and relation
\[
 U R^{\mathrm{ue}}U'\quad\text{iff}\quad
 \forall X\subseteq W\,(l(X)\in U\Rightarrow X\in U').
\]
For a model on \(F\), its extension gives an atom \(p\) the set
\(\{U:V(p)\in U\}\). Ultrafilter existence uses the Boolean
prime ideal theorem, available in the standing ZFC metatheory.
\end{hbdefinition}

\begin{hbtheorem}{thm:p08-bvb501-ultrafilter-reflection}{proved=annex:p08.src025501.ultrafilter}{Source Theorem 30 and the natural-number example}
For every formula, truth in an extended model at \(U\) means its
original truth set belongs to \(U\). In particular, truth at the
principal ultrafilter of \(w\) agrees with original truth at \(w\).
Thus \(\operatorname{ue}(F)\entails\varphi\) implies
\(F\entails\varphi\): validity reflects ultrafilter extensions.
The class \(\forall x\exists y(Rxy\land Ryy)\) fails this
reflection and is not modally definable.
\end{hbtheorem}

The source's ``lolly-pop'' picture for \(\operatorname{ue}(\mathbb N,<)\)
has principal points in their original order, followed by nonprincipal
ultrafilters forming a reflexive universal cluster. Every principal point
accesses every member of that cluster. The original order has no loop,
but every extension point has a reflexive successor. The annex verifies
the relation, including the existence of many nonprincipal points.

\subsection{Second-order and first-order correspondence}
\label{sec:p08-bvb501-correspondence}
\index{correspondence language!second-order}
\index{frame correspondence language}
\index{Loeb's formula@L\"ob's formula}
\index{McKinsey formula}
\index{Sahlqvist Correspondence Theorem}

\begin{hbtheorem}{thm:p08-bvb501-frame-second-order}{proved=annex:p08.src025501.frame-so}{The universal monadic second-order closure}
If \(p_1,\ldots,p_n\) are the formula's proposition letters, then
\[
 F\entails\varphi\quad\text{iff}\quad
 F\entails\forall P_1\cdots\forall P_n\forall x\,
       \operatorname{ST}_x(\varphi).
\]
The predicate quantifiers range over all subsets of the frame domain.
For \(p\to\Diamond p\), choosing the singleton valuation
\(P=\{x\}\) reduces this second-order condition to reflexivity.
\end{hbtheorem}

\begin{hbtheorem}{thm:p08-bvb501-loeb-correspondence}{proved=annex:p08.src025501.loeb-frame}{L\"ob's nonelementary frame condition}
\(\Box(\Box p\to p)\to\Box p\) defines transitive frames
with no infinite forward \(R\)-chain. This class is not elementary,
even though the modal formula is finite. The source separately cites
McKinsey's \(\Box\Diamond p\to\Diamond\Box p\) as another
formula with no first-order frame correspondent on arbitrary frames.
\end{hbtheorem}

\begin{hbtheorem}{thm:p08-bvb501-sahlqvist-correspondence}{cited=BlackburnVanBenthem2006ILLC}{Source Theorem 31 and its stated syntactic class}
An implication whose antecedent is built from boxed atoms by conjunction,
disjunction and diamonds, and whose consequent has only positive atom
occurrences, has an effectively computable first-order frame correspondent.
The source's substitution method uses minimal antecedent valuations
after fixing diamond witnesses and disjunctive branches.
\end{hbtheorem}

For \(\Box p\to\Box\Box p\) at \(x\), the minimal
antecedent valuation is \(P(u)\) iff \(Rxu\). Substitution in
\(\forall P\forall x(\forall y(Rxy\to Py)\to
\forall y(Rxy\to\forall z(Ryz\to Pz)))\) leaves
\(\forall xyz(Rxy\land Ryz\to Rxz)\). The complete necessity
and sufficiency are proved in the correspondent theorem above. Arbitrary
disjunctions and diamonds need witness-dependent branches; they need
not have a single least antecedent valuation before witnesses are fixed.

\begin{hbtheorem}{thm:p08-bvb501-first-order-boundaries}{cited=BlackburnVanBenthem2006ILLC}{Source Theorems 32--33 and the non-Sahlqvist example}
A modal formula's valid-frame class is elementary iff it is closed
under ultrapowers. Determining whether a modal formula has a first-order
frame correspondent is undecidable. For an elementary frame class,
modal definability is equivalent to closure under disjoint unions,
generated subframes and bounded morphic images, and reflection of
ultrafilter extensions (Goldblatt--Thomason).

The conjunction of transitivity and McKinsey's axiom has a first-order
correspondent despite lying outside the displayed Sahlqvist class:
\(R\) is transitive and
\(\forall x\exists y(Rxy\land\forall z(Ryz\to z=y))\).
The source calls this the transitive atomic condition and refers its
proof to van Benthem and Blackburn, de Rijke and Venema.
\end{hbtheorem}

The necessity of the Goldblatt--Thomason closures is proved here.
The sufficiency is the source's cited algebraic or model-theoretic
representation theorem, not a consequence of the closure examples alone.
Likewise, the ultrapower criterion is not an algorithm for recognizing
first-order correspondents.

\subsection{Fixed-point correspondents and computability}
\label{sec:p08-bvb501-fixed-correspondence}
\index{global frame consequence}

\begin{hbdefinition}{def:p08-bvb501-intersective}{Intersection-preserving antecedent conditions}
For fixed parameters \(Q\), a closure condition
\(\forall x(\chi(P,Q,x)\to Px)\), with \(P\) positive
in \(\chi\), is preserved under nonempty intersections of
satisfying predicates. A condition with only negative occurrences of
\(P\) is likewise downward closed. A least satisfying predicate
exists in the first case because \(W\) satisfies the closure condition;
in the negative case existence additionally requires some satisfying
predicate. An inconsistent negative condition has none.
\end{hbdefinition}

\begin{hbtheorem}{thm:p08-bvb501-minimal-fixed-point}{proved=annex:p08.src025501.minimum}{Minimal closure and cyclic-return correspondence}
For a monotone operator \(G\) on \(\mathcal P(W)\), the
intersection of all \(P\) with \(G(P)\subseteq P\) is its
least fixed point. Hence the positive closure conditions above have
least predicates definable by a least-fixed-point operator. The axiom
\((\Diamond p\land\Box(p\to\Box p))\to p\) defines
\(\forall xy(Rxy\to R_x^*yx)\), where
\(R_xuz\) means \(Rxu\land Ruz\): every successor can return
to its source along a finite path whose departures remain successors
of that source. On transitive frames this is equivalent to ordinary
finite return \(\forall xy(Rxy\to R^*yx)\).
\end{hbtheorem}

For cyclic return, first fix a diamond witness \(y\). Its least
predicate closed under departures among \(x\)'s successors is
the set reachable from \(y\) by \(R_x\). The
whole antecedent with an unfixed existential witness is not itself
intersection preserving: a root with two separate loop successors
has two disjoint satisfying predicates, while their intersection fails
the diamond. Pulling out the witness is an essential part of the
minimal-predicate computation. The ordinary-return simplification needs
transitivity: on the directed cycle \(xRy,yRz,zRx\) every edge
returns, but the valuation \(\{y,z\}\) makes the axiom false at
\(x\). The \(R_x\)-path from \(y\) cannot continue through
\(z\), which is not a successor of \(x\). For L\"ob's antecedent at \(x\),
the closure operator adds those \(R\)-successors of \(x\) whose
successors all belong to \(P\). Its least fixed point supplies a
correspondent in first-order least-fixed-point logic rather than ordinary
first-order logic. The source cites extensions of this analysis and
the failure of an LFP(FO) correspondent for McKinsey. It also contrasts
Dedekind completeness of linear orders: rational and real orders are
potentially isomorphic but disagree on completeness, so potential-
isomorphism-invariant LFP(FO) cannot define that condition.

\begin{hbtheorem}{thm:p08-bvb501-functional-complexity}{proved=annex:p08.src025501.functional}{Partial-function frames give NP-complete satisfiability}
For one basic modality constrained to a partial-function relation,
every satisfiable formula has a path model with at most
\(\operatorname{md}(\varphi)+1\) states. Satisfiability on this
class is NP-complete. The depth counts nested modal operators, not the
number of distinct modality names.
\end{hbtheorem}

\begin{hbtheorem}{thm:p08-bvb501-elementary-enumeration}{proved=annex:p08.src025501.frame-enumerate}{Source Proposition 34 with its effectiveness condition}
For an effective basic modal signature, if an elementary frame class is
specified by a first-order sentence
\(\alpha\), its basic modal validities are recursively enumerable.
The same conclusion holds for a computably axiomatized first-order
frame theory. Bare elementarity without an effective specification
does not supply an enumeration procedure.
\end{hbtheorem}

Restricting a frame class can lower complexity, raise it, or make the
decision problem undecidable, including for some finite-frame classes.
The source reports uncountably many distinct frame satisfiability
problems; only countably many algorithms exist. This cardinality
observation concerns a family of classes, not uncountably many individual
finite modal formulas. Beyond elementary classes, Thomason's reduction
of second-order consequence to global basic-modal frame consequence
shows that even recursive enumerability can fail. Global frame
consequence quantifies over frames that validate the premises under
every valuation; it is not the local model consequence relation.

\section{Richer languages and their design choices}
\label{sec:p08-bvb501-richer-languages}
\index{polyadic modalities}

Richer logics add axioms in a fixed language; richer languages change
what can be expressed. The source separates fixed interpretations of
special relations, naming devices, temporal interval operators,
conditionals, guarded quantification, programs, fixed points, combinations
and first-order modal quantification. They need not have the same
decidability or invariance properties.

\begin{hbdefinition}{def:p08-bvb501-polyadic}{The preliminary polyadic variant}
An \(n\)-place diamond uses an \((n+1)\)-place relation:
\[
 M,w\entails\langle m\rangle(\varphi_1,\ldots,\varphi_n)
 \quad\text{iff}\quad
 \exists v_1\cdots v_n\,(R_mwv_1\cdots v_n\land
             \bigwedge_{i=1}^n M,v_i\entails\varphi_i).
\]
Its standard translation quantifies that successor tuple and translates
each argument at its corresponding coordinate. A bisimulation matches
each successor tuple with a tuple whose corresponding coordinates are
related, in both directions. Atomic harmony and these tuple clauses give
truth invariance by the same structural induction as for unary diamonds.
This changes arity without yet introducing naming, path iteration or
fixed-point expressivity.
\end{hbdefinition}

\subsection{The universal and difference modalities}
\label{sec:p08-bvb501-universal}
\index{universal modality}
\index{difference operator}

\begin{hbdefinition}{def:p08-bvb501-universal}{Global truth as an object-language modality}
Add \(E\) and its dual \(A\), always interpreted by \(W^2\):
\[
 M,w\entails E\varphi\quad\text{iff}\quad
 \exists u\in W\ M,u\entails\varphi,\qquad
 M,w\entails A\varphi\quad\text{iff}\quad
 \forall u\in W\ M,u\entails\varphi.
\]
The difference diamond \(D\) instead uses inequality:
\(M,w\entails D\varphi\) iff some \(u\ne w\) satisfies
\(\varphi\). Neither interpretation is a freely chosen extra
accessibility relation. A total bisimulation covers every point of
each model by some related point in the other; it need not be functional
or injective.
\end{hbdefinition}

\begin{hbtheorem}{thm:p08-bvb501-universal-expressivity}{proved=annex:p08.src025501.universal}{Strict enrichment and total-bisimulation invariance}
The universal modality is not definable in the basic language. The
enriched language is invariant under total bisimulations. Local--global
constraints \(\varphi\) locally and \(\psi\) globally are
expressed by \(\varphi\land A\psi\); global satisfiability of
\(\psi\) is satisfiability of \(A\psi\).
\end{hbtheorem}

This internalizes the background-constraint layer used in the source's
BEAR-CARE example. A fresh inaccessible component can change \(E\)
or \(A\) truth, although it cannot change basic modal truth at a
point of the original component. Expressivity therefore increases
without any assumption separating complexity classes.

\begin{hbtheorem}{thm:p08-bvb501-universal-complexity}{cited=BlackburnVanBenthem2006ILLC}{Source Theorem 35}
Satisfiability for the basic language with the universal modality is
EXPTIME-complete. Its lower bound includes encoding the global and
local--global constraints; the source refers the matching upper bound
and full classification to Spaan and Hemaspaandra.
\end{hbtheorem}

\begin{hbtheorem}{thm:p08-bvb501-universal-frames}{proved=annex:p08.src025501.universal-frames}{The source's newly definable cardinality and predecessor conditions}
For \(n\geq1\), the formula
\[
 \left(\bigwedge_{i=1}^{n+1}Ep_i\right)
       \to\bigvee_{1\leq i<j\leq n+1}E(p_i\land p_j)
\]
defines frames with at most \(n\) states. The formula
\(p\to E\Diamond p\) defines
\(\forall x\exists y\,Ryx\), existence of a predecessor
for every point. Neither class is basic modally definable: the first
fails closure under disjoint unions and the second fails closure under
generated subframes. These are failures of old preservation conditions,
not failures of the new interpretation.
\end{hbtheorem}

For \(n=0\), the bounded-cardinality class is empty under the standing
nonempty-frame convention, and is defined by \(\bottom\).

\begin{hbtheorem}{thm:p08-bvb501-universal-gt}{cited=BlackburnVanBenthem2006ILLC}{Source Theorem 36}
An elementary frame class is definable using the universal modality iff
it is closed under surjective bounded morphic images and reflects
ultrafilter extensions. The source attributes this variant to Goranko
and Passy. Disjoint-union and generated-subframe closure are absent.
\end{hbtheorem}

\begin{hbtheorem}{thm:p08-bvb501-universal-well-founded}{proved=annex:p08.src025501.universal-wf}{Well-foundedness without requiring transitivity}
The formula \(A(\Box p\to p)\to p\) defines frames with no
infinite forward \(R\)-chain, without adding transitivity. That class
is not definable in the basic language. Its basic validities are exactly
K's validities, since every falsifiable basic formula has a finite
acyclic rooted tree countermodel.
\end{hbtheorem}

\begin{hbtheorem}{thm:p08-bvb501-difference}{proved=annex:p08.src025501.difference}{Difference defines the universal modality, but not conversely}
\(E\varphi\leftrightarrow(\varphi\lor D\varphi)\) is
valid. There is no definition of \(D\) using just basic modalities
and \(E,A\). In particular, \(D\topv\) distinguishes a
singleton model from a two-point model with otherwise identical isolated
states; a total bisimulation preserves all formulas of the smaller language.
\end{hbtheorem}

\subsection{Hybrid naming and equality}
\label{sec:p08-bvb501-hybrid}
\index{hybrid logic}
\index{basic hybrid language}
\index{nominals}
\index{satisfaction operators}
\index{Hybrid Characterisation Theorem}

\begin{hbdefinition}{def:p08-bvb501-hybrid}{Nominals, satisfaction operators and named bisimulations}
Add nominals \(i,j,k\), each assigned a singleton \(\{d_i\}\),
and operators \(@_i\), with \(M,w\entails @_i\varphi\) iff
\(M,d_i\entails\varphi\). Different nominals may name the same
point. A bisimulation with names checks both ordinary atoms and nominals
at related points and additionally relates \(d_i\) to \(d_i'\)
for every nominal, even if these points are unreachable from the selected
point. In the correspondence language, treat \(i\) as a constant:
\(\operatorname{ST}_x(i)=(x=i)\) and
\(\operatorname{ST}_x(@_i\varphi)=\operatorname{ST}_i(\varphi)\),
with capture-free substitution for the displayed free variable.
\end{hbdefinition}

\begin{hbtheorem}{thm:p08-bvb501-hybrid-basic}{proved=annex:p08.src025501.hybrid}{Source Lemma 38 and the elementary naming consequences}
The extended translation preserves truth and has at most one free world
variable. Bisimulations with names preserve every basic hybrid formula.
Satisfaction operators distribute over implication, admit theorem-only
necessitation, and are self-dual. The formulas \(@_i i\),
\(@_i j\to @_j i\), \((@_i j\land @_j k)\to @_i k\),
and \((@_i\varphi\land @_i j)\to @_j\varphi\) are valid.
The formula \(\Diamond(i\land p)\land\Diamond(i\land q)
\to\Diamond(p\land q)\) is valid, unlike its ordinary-atom version.
\end{hbtheorem}

Naming also expands frame definability. The respective hybrid formulas
for irreflexivity, asymmetry, antisymmetry and trichotomy are
\(i\to\neg\Diamond i\), \(i\to\neg\Diamond\Diamond i\),
\(i\to\Box(\Diamond i\to i)\), and
\(@_j\Diamond i\lor @_j i\lor @_i\Diamond j\).
Frame validity now quantifies over every assignment of names to points.

\begin{hbtheorem}{thm:p08-bvb501-hybrid-frames}{proved=annex:p08.src025501.hybrid-frames}{The four named frame correspondents}
These formulas define, respectively,
\(\forall x\neg Rxx\), \(\forall xy(Rxy\to\neg Ryx)\),
\(\forall xy(Rxy\land Ryx\to x=y)\), and
\(\forall xy(Rxy\lor x=y\lor Ryx)\).
\end{hbtheorem}

\begin{hbtheorem}{thm:p08-bvb501-hybrid-characterization}{cited=BlackburnVanBenthem2006ILLC}{Source Theorems 39--40}
Basic hybrid formulas characterize exactly the fragment of first-order
correspondence logic with constants and equality invariant under
bisimulations with names. Satisfiability over arbitrary hybrid models is
PSPACE-complete. The source attributes both results to Areces, Blackburn
and Marx; neither says that arbitrary first-order equality reasoning has
become a decidable hybrid satisfiability problem.
\end{hbtheorem}

\subsection{Temporal Until and Since}
\label{sec:p08-bvb501-temporal}
\index{Dedekind completeness}
\index{expressive completeness}
\index{Kamp's theorem}

\begin{hbdefinition}{def:p08-bvb501-until-since}{The source's strict interval operators}
On an arbitrary relation \(R\), let
\[
 \begin{aligned}
 M,w\entails U(\varphi,\psi)&\quad\text{iff}\quad
  \exists v(wRv\land M,v\entails\varphi\land
     \forall u(wRu\land uRv\Rightarrow M,u\entails\psi)),\\
 M,w\entails S(\varphi,\psi)&\quad\text{iff}\quad
  \exists v(vRw\land M,v\entails\varphi\land
     \forall u(vRu\land uRw\Rightarrow M,u\entails\psi)).
 \end{aligned}
\]
These inspect the open relational interval, not the existential finite
path of the queue algorithm. On a strict temporal order \(U\)
means a future endpoint satisfying \(\varphi\), with \(\psi\)
at every strictly intermediate point; \(S\) reverses time. The
order need not have an immediate next point. \(U(\varphi,\topv)\)
is the ordinary future diamond and \(S(\varphi,\topv)\) the past
diamond. A Dedekind-complete strict total order has a supremum for every
nonempty subset bounded above; excluding the empty subset is necessary
for the real and natural-number examples.
\end{hbdefinition}

\begin{hbtheorem}{thm:p08-bvb501-temporal-translation}{proved=annex:p08.src025501.temporal}{Three-variable translation and strict enrichment}
The truth-preserving clauses are
\[
 \begin{aligned}
 \operatorname{ST}_x(U(\varphi,\psi))&=
  \exists z(Rxz\land\operatorname{ST}_z(\varphi)\land
       \forall y(Rxy\land Ryz\to\operatorname{ST}_y(\psi))),\\
 \operatorname{ST}_x(S(\varphi,\psi))&=
  \exists z(Rzx\land\operatorname{ST}_z(\varphi)\land
       \forall y(Rzy\land Ryx\to\operatorname{ST}_y(\psi))).
 \end{aligned}
\]
Bound variables are renamed locally to avoid capture. Until is not
definable in the basic future modal language, even on finite transitive
irreflexive frames.
\end{hbtheorem}

For the separating construction based on source Figure~16, the left graph has arrows
\(w_0\to t_0,u\), \(w_1\to u,t_1\),
\(t_0\to v_0\), \(u\to v_0,v_1\), \(t_1\to v_1\),
with their transitive closure. Only \(v_0,v_1\) satisfy \(p\);
only \(u\) satisfies \(q\). On the right, \(w'\) leads to
\(u',t'\), each leading to \(v'\); only \(v'\) satisfies
\(p\), and only \(u'\) satisfies \(q\). Relate both roots to
\(w'\), \(u\) to \(u'\), \(t_0,t_1\) to \(t'\)
and both endpoints to \(v'\). This is a future bisimulation.
At \(w_0\), choose endpoint \(v_1\): its only intermediate
point is \(u\), satisfying \(q\). At \(w_1\), choose
\(v_0\), again with only \(u\) intermediate. At \(w'\) the only
\(p\)-endpoint has the false-\(q\) intermediate \(t'\).
Thus \(U(p,q)\) separates the related roots. The source additionally
cites undefinability even on the real order.

\begin{hbtheorem}{thm:p08-bvb501-kamp}{cited=BlackburnVanBenthem2006ILLC}{Source Theorem 41 and its specified order class}
Basic temporal logic enriched by Until and Since is expressively complete
for the one-free-variable first-order correspondence language over
Dedekind-complete strict total orders, with arbitrary unary valuations.
The source refers to Kamp's 1968 thesis and Gabbay, Hodkinson and Reynolds.
It also cites decidability over arbitrary relational models and
Until/Since comparison notions due to Kurtonina--de Rijke and Sturm.
\end{hbtheorem}

The source distinguishes expressivity from the use of these operators
for program guarantees, and names LTL, CTL and CTL* as temporal languages
with related but different path/state semantics. Kamp's theorem is not
a statement identifying these three computational languages.

\subsection{Conditional preference and belief revision}
\label{sec:p08-bvb501-conditional}
\index{conditional logic}

\begin{hbdefinition}{def:p08-bvb501-conditional}{Closer worlds and the no-limit truth clause}
Let \(C_wuv\) mean that \(u\) is at least as close or preferred
to the reference world \(w\) as \(v\) is. Each \(C_w\)
is a reflexive transitive preorder; comparability is an additional
condition. A minimal antecedent world has no strictly better antecedent
world, where strict preference means \(C_wuv\land\neg C_wvu\).
On finite preorders, \(\varphi>\psi\) means that all minimal
\(\varphi\)-worlds satisfy \(\psi\). For arbitrary preorders
the threshold clause is
\[
 \forall u\bigl(\varphi(u)\Rightarrow
   \exists v(C_wvu\land\varphi(v)\land
     \forall z(C_wzv\land\varphi(z)\Rightarrow\psi(z)))\bigr),
\]
where \(\varphi(t)\) denotes truth in the original conditional
model at \(t\). Both the closer-to-\(v\) relation and the
antecedent test concern \(z\) in the innermost clause. Empty
antecedent extensions make the conditional vacuously true.
\end{hbdefinition}

\begin{hbtheorem}{thm:p08-bvb501-conditional-clause}{proved=annex:p08.src025501.conditional}{Thresholds, finite minima and the correct centered modal reading}
On finite preorders the threshold clause is equivalent to truth of
\(\psi\) at every minimal antecedent world. It is upward monotone
in the consequent but not in the antecedent. For a fixed reference
\(w\), put \(uR_w^cv\) iff \(C_wvu\). The clause then
has the auxiliary modal reading
\(A(\varphi\to\langle c\rangle(\varphi\land[c](\varphi\to\psi)))\),
with the already computed original argument truth sets. Nested conditional
evaluation changes the reference world and is not obtained by globally
fixing this one auxiliary relation.
\end{hbtheorem}

The source's Amsterdam-train and broken-dike example illustrates this
failure of antecedent strengthening. An infinite improving chain may
have no minimal antecedent world; the threshold clause handles that
case without declaring every consequent true merely because no minimum
exists. The source discusses Burgess and Veltman's minimal systems,
Lewis's comparable similarity orders, geometrical correspondents, and
belief-revision readings due to G\"ardenfors--Rott and Ryan--Schobbens.
Different order/domain restrictions specify different conditional logics.
In variants allowing a world-specific possible set, the preference
preorder and its conditional quantifiers are restricted to that set.
A preorder on all of the common \(W\) is already a possible-domain
restriction; varying-domain complexity claims cannot simply be inferred
for it by forgetting that distinction.

\begin{hbtheorem}{thm:p08-bvb501-conditional-complexity}{cited=BlackburnVanBenthem2006ILLC}{Source Theorem 42: the stated conditional-complexity variants}
The source reports PSPACE-completeness with unrestricted conditional
nesting and NP-completeness with a fixed nesting bound for the minimal
ordering variants; uniform possible-world domains yield EXPTIME-complete
variants even at bounded nesting, while absolute conditional structures
yield NP-completeness with arbitrary nesting. It cites Friedman and Halpern
for the exact structures and qualifications. Uniformity concerns agreement
on which worlds are possible, and absoluteness agreement on conditional
judgments; neither is mere preorder transitivity or comparability.
\end{hbtheorem}

\subsection{Guarded and packed first-order quantification}
\label{sec:p08-bvb501-guarded}
\index{guarded fragment}
\index{packed fragment}

\begin{hbdefinition}{def:p08-bvb501-guarded}{The two source restricted-quantification fragments}
Guarded formulas begin with atomic formulas and equality, close under
Booleans, and allow \(\exists\bar y(G(\bar x,\bar y)\land\chi)\)
and \(\forall\bar y(G(\bar x,\bar y)\to\chi)\).
The atomic guard contains every free variable of the matrix and every
quantified variable retained in the quantification; \(\bar x,\bar y\)
are finite tuples. Packed guards instead are conjunctions of atoms,
equalities, or quantified relational atoms, with the additional condition
that each pair of distinct free variables in the guard occurs together
free in some conjunct. The guard still covers the matrix's free variables.
This pair-covering requirement is stronger than merely putting each
variable somewhere in a conjunction.
\end{hbdefinition}

\begin{hbtheorem}{thm:p08-bvb501-packed-example}{proved=annex:p08.src025501.packed}{The corrected Until packed translation}
The Until translation
\(\exists y(Rxy\land Py\land\forall z(Rxz\land Rzy\to Qz))\)
is equivalent to
\(\exists y(Rxy\land Py\land
\forall z(Rxz\land Rzy\land Rxy\to Qz))\).
The latter is packed: the inner guard covers pairs \(x,z\),
\(z,y\) and \(x,y\). The outer quantifier binds \(y\),
not \(x\), so the formula retains its intended free evaluation variable.
The original inner guard is not atomic and is not packed.
\end{hbtheorem}

\begin{hbtheorem}{thm:p08-bvb501-guarded-results}{cited=BlackburnVanBenthem2006ILLC}{Source Theorems 43--44 and the metaproperty boundaries}
The guarded fragment has the finite model property and 2EXPTIME-complete
satisfiability; a fixed predicate-arity bound gives EXPTIME-completeness.
The packed fragment has the finite model property and 2EXPTIME-complete
satisfiability. The source attributes these results to Gr\"adel,
Andr\'eka--Hodkinson--N\'emeti and Hodkinson--Otto.
It cites guarded-bisimulation characterization and Beth definability,
but distinguishes failure of unrestricted Craig interpolation from
interpolation with guards held as logical vocabulary. The inability
to express Until in the guarded fragment is an expressivity result,
not just the failure of its displayed formula to match the grammar.
\end{hbtheorem}

\subsection{Propositional dynamic logic}
\label{sec:p08-bvb501-pdl}
\index{propositional dynamic logic}
\index{Segerberg's axiom}

\begin{hbdefinition}{def:p08-bvb501-pdl}{Programs, tests and regular models}
Define programs and formulas mutually. Atomic programs \(a\) are
closed under choice \(\pi\cup\rho\), composition \(\pi;\rho\),
finite iteration \(\pi^*\), and tests \(\chi?\) for formulas
\(\chi\). Formulas use Booleans and \(\langle\pi\rangle\chi\).
An atomic program has an arbitrary binary transition relation; composite
relations are
\[
 R_{\pi\cup\rho}=R_\pi\cup R_\rho,\quad
 R_{\pi;\rho}=R_\pi\circ R_\rho,\quad
 R_{\pi^*}=(R_\pi)^*,\quad
 R_{\chi?}=\{(w,w):M,w\entails\chi\}.
\]
Here composition executes its left program first. Tests depend on the
valuation and are not independent relations of a valuation-free frame.
For test-free programs these equations define a regular frame; with
tests they define a regular model by mutual semantic recursion.
Diamonds assert some terminating execution, boxes every terminating
execution; neither asserts that an execution terminates.
\end{hbdefinition}

The source's program encodings are
\((p?;a)\cup(\neg p?;b)\) for an if--then--else,
\(a;(\neg p?;a)^*\) for repeat-until (stopping after a
\(p\)-test may be imposed explicitly if termination at \(p\)
is intended), and \((p?;a)^*;\neg p?\) for while-\(p\)-do-\(a\).
The untested repeat expression permits premature stopping after an
\(a\)-execution; the conventional terminating repeat-until program
is \(a;(\neg p?;a)^*;p?\). Star denotes finite iteration, including
zero, not an infinite run.

\begin{hbtheorem}{thm:p08-bvb501-pdl-regularity}{proved=annex:p08.src025501.pdl-regular}{The source regularity axioms and their exact force}
For independently named test-free program relations, frame validity of
the choice and composition diamond equations forces the respective
relational union and composition. The star unfolding equation
\(\langle\pi^*\rangle p\leftrightarrow
(p\lor\langle\pi;\pi^*\rangle p)\), together with composition
and Segerberg's induction
\([\pi^*](p\to[\pi]p)\to(p\to[\pi^*]p)\), forces
\(R_{\pi^*}=(R_\pi)^*\). Thus those schemas characterize
regular test-free frames. Arbitrary transitive closure is not uniformly
first-order definable.
\end{hbtheorem}

\begin{hbtheorem}{thm:p08-bvb501-pdl-results}{cited=BlackburnVanBenthem2006ILLC}{Source Theorem 45}
PDL has the finite model property and EXPTIME-complete satisfiability.
The source attributes the finite-model and lower-bound results to
Fischer and Ladner and the upper bound to Pratt. The test constructor
and regular-model conditions are part of this specified language.
\end{hbtheorem}

\subsection{The modal fixed-point calculus}
\label{sec:p08-bvb501-mu}
\index{modal mu-calculus}
\index{Knaster--Tarski theorem}
\index{positive occurrences}
\index{negative occurrences}

\begin{hbdefinition}{def:p08-bvb501-mu}{Positive binders and least/greatest solutions}
For a formula \(\chi\) and proposition variable \(p\), the
operator it induces in a fixed model is
\(F(U)=V[p\leftarrow U](\chi)\), where the extended valuation
returns truth sets. Allow \(\mu p.\chi\) only when every free
occurrence of \(p\) in \(\chi\) is positive; expand redundant
Boolean constructors before computing polarity. Bound variables are
renamed to prevent capture. Then
\[
 V(\mu p.\chi)=\bigcap\{U:F(U)\subseteq U\},\qquad
 V(\nu p.\chi)=\bigcup\{U:U\subseteq F(U)\}.
\]
The greatest binder is equivalently
\(\nu p.\chi=\neg\mu p.\neg\chi[\neg p/p]\), with
substitution restricted to free occurrences. Negative dependence need
not have any fixed point: complement on a nonempty powerset has none.
Write \(\operatorname{pre}_a(X)=\{w:\exists v(wR_av\land v\in X)\}\)
and \(\operatorname{box}_a(X)=\{w:\forall v(wR_av\Rightarrow v\in X)\}\)
for the existential and universal successor-set operations.
\end{hbdefinition}

\begin{hbtheorem}{thm:p08-bvb501-mu-constructions}{proved=annex:p08.src025501.mu}{Fixed-point semantics, PDL translation and the MSO clause}
Positivity makes the induced operator monotone; the displayed extrema
are its least and greatest fixed points. With fresh \(p\),
\(\mu p.(\chi\lor\langle a\rangle p)\) denotes finite
reachability of \(\chi\), \(\nu p.(\chi\land[a]p)\)
denotes \(\chi\) at every finitely reachable point, and
\(\nu p.\langle a\rangle p\) denotes existence of an infinite
\(a\)-run. The calculus is bisimulation invariant.

Translate PDL mutually through a continuation operation \(T_\pi(\eta)\):
\[
 \begin{aligned}
 T_a(\eta)&=\langle a\rangle\eta,&
 T_{\pi\cup\rho}(\eta)&=T_\pi(\eta)\lor T_\rho(\eta),\\
 T_{\pi;\rho}(\eta)&=T_\pi(T_\rho(\eta)),&
 T_{\chi?}(\eta)&=T(\chi)\land\eta,\\
 T_{\pi^*}(\eta)&=\mu p.(\eta\lor T_\pi(p)),&
 T(\langle\pi\rangle\chi)&=T_\pi(T(\chi)).
 \end{aligned}
\]
Here \(p\) is fresh in the continuation, programs and their tests;
the Boolean clauses of \(T\) commute with their constructors.
This translation preserves truth and eliminates composite program
modalities and tests from the target.

The monadic second-order standard translation for a least binder is
\[
 \operatorname{ST}_x(\mu p.\chi)=
 \forall P\left((\forall y(\operatorname{ST}_y(\chi)\to Py))\to Px\right),
\]
where free occurrences of \(p\) translate to the fresh bound predicate
\(P\), and other variables/predicates are renamed hygienically.
It asserts membership in every pre-fixed predicate, not merely an
implication evaluated at the original free variable.
\end{hbtheorem}

\begin{hbtheorem}{thm:p08-bvb501-mu-results}{cited=BlackburnVanBenthem2006ILLC}{Source Theorems 46--47 and the stated expressivity comparisons}
The modal \(\mu\)-calculus has the finite model property and
EXPTIME-complete satisfiability. It is the bisimulation-invariant
fragment of monadic second-order logic on pointed relational models.
PDL is strictly weaker; the source cites infinite-run repetition as
a separating property and embeddings of computational temporal logics.
It attributes these results to Kozen--Parikh, Streett--Emerson,
Emerson--Jutla and Janin--Walukiewicz. Its 2006 discussion of a
polynomial model-checking conjecture concerns a different problem from
the stated satisfiability classification and is historical, not an
undated current complexity assertion.
\end{hbtheorem}

\subsection{Fusions, interactions and products}
\label{sec:p08-bvb501-combinations}

\begin{hbdefinition}{def:p08-bvb501-fusion-product}{Two different combination operations}
For normal logics with disjoint modality alphabets, their fusion
\(L_1\otimes L_2\) is the least normal logic in the joint language
containing both, with uniform substitution by joint-language formulas.
It introduces no additional mixed axioms. A product of frames instead
has domain \(W_1\times W_2\) and
\[
 (u,v)R_h(u',v')\ \text{iff}\ uR_1u'\land v=v',\qquad
 (u,v)R_v(u',v')\ \text{iff}\ vR_2v'\land u=u'.
\]
Horizontal and vertical steps commute. This strong interaction is not
part of the definition of fusion.
\end{hbdefinition}

The source's planning examples contrast \(K_i[a]\varphi\), knowing
an action's effect, with \([a]K_i\varphi\), gaining knowledge after
acting. Neither implies the other on arbitrary combined frames.

\begin{hbtheorem}{thm:p08-bvb501-interactions}{proved=annex:p08.src025501.interactions}{The two mixed frame correspondents and counterexamples}
Writing \(I\) for agent \(i\)'s relation and \(R_a\) for
the action, frame validity of \(K_i[a]p\to[a]K_ip\) is exactly
\(R_a\circ I\subseteq I\circ R_a\); validity of
\([a]K_ip\to K_i[a]p\) is the reverse inclusion. Both can fail
even when the epistemic relation is an equivalence and the action a
partial function. Product frames satisfy both commuting inclusions.
\end{hbtheorem}

The source associates these patterns with perfect recall and no learning.
For its interpreted-system settings it cites Halpern--Vardi's
nonelementary complexity with perfect recall and undecidability after
adding common knowledge. Common knowledge uses finite closure of the
union of agents' relations. Those are results for specified interaction
classes, not a theorem that every fusion is undecidable. For standard
normal fusions the source reports transfer of decidability and the
finite model property, citing Kracht--Wolter, Fine--Schurz and Wolter.
Finite axiomatizations combine by taking their union. Products and
intermediate interaction methods need separate transfer arguments;
grid-like interactions can support tiling encodings. The source credits
Moore's knowledge/action language and Segerberg's early products.

\subsection{Quantification inside modal contexts}
\label{sec:p08-bvb501-first-order-modal}
\index{first-order modal logic}

\begin{hbdefinition}{def:p08-bvb501-varying-domain}{The source's outer-domain varying-domain semantics}
Use \(M=(W,R,D,(\delta_w),(V_w))\), with nonempty \(W,D\),
\(\delta_w\subseteq D\), and \(V_w(P)\subseteq D^n\)
for each \(n\)-ary predicate. The source convention permits empty
local domains; a nonempty-local-domain version is a separate restriction.
Assignments \(g\) take values in \(D\), remain unchanged across
modal moves, and equality compares their values. Atoms use \(V_w\),
Booleans are classical, and modal clauses quantify over accessible
worlds. Quantifiers stay at the current world:
\[
 \begin{aligned}
 M,g,w\entails\exists x\chi&\ \text{iff}\ 
   \exists d\in\delta_w\ M,g[x\mapsto d],w\entails\chi,\\
 M,g,w\entails\forall x\chi&\ \text{iff}\ 
   \forall d\in\delta_w\ M,g[x\mapsto d],w\entails\chi.
 \end{aligned}
\]
Thus binding can reach occurrences inside modal scopes. Predicate
values on all of \(D^n\) are defined, even when some assigned
individual is absent from a world's quantifier domain.
\end{hbdefinition}

\begin{hbtheorem}{thm:p08-bvb501-barcan}{proved=annex:p08.src025501.barcan}{Source Barcan correspondences with the valuation quantifier explicit}
Over a fixed world relation and domain family, all instances of
\(\forall x\Box\chi\to\Box\forall x\chi\) hold under
every predicate interpretation and assignment iff accessible domains
decrease: \(wRv\Rightarrow\delta_v\subseteq\delta_w\).
All converse instances \(\Box\forall x\chi\to\forall x\Box\chi\)
hold iff accessible domains increase. Both force locally constant domains,
not necessarily equal domains in disconnected components.
\end{hbtheorem}

The source's distinction is between statements about each currently
quantified individual across worlds and statements quantifying anew at
each accessible world. Its philosophical background includes Quine's
objections to binding into intensional contexts and Kripke's relational
semantics and \emph{Naming and Necessity}. It cites Fine for failures of
interpolation/Beth properties in varying-domain quantified S5 and
constant-domain extensions between K and S5, and Ghilardi,
Shehtman--Skvortsov, Corsi--Ghilardi and Cresswell for failure of
completeness transfer. Propositional completeness does not automatically
survive this combination with quantification. Functional-frame,
categorical and first-order hybrid approaches are alternative settings;
the source credits van Benthem, Areces--Blackburn--Marx and
Blackburn--Marx for particular repairs, without asserting that every
first-order modal logic thereby gains interpolation or completeness.

\subsection{Guarding and abstract modal characterizations}
\label{sec:p08-bvb501-abstract}

Finite-variable translation, bisimulation unraveling and guarding are
different explanatory tools. A fixed number of variables supports
polynomial combined model checking by maintaining truth tables for
assignments, but three-variable first-order satisfiability is already
undecidable. Variable count alone therefore does not explain modal
decidability. The source cites Vardi's tree-based explanation and
Gr\"adel--Walukiewicz's decidable guarded fixed-point logic, contrasting
full first-order fixed-point validity's nonarithmetical behavior. Guarding
is a useful design heuristic, not a proved definition covering every
language called modal. Infinitary languages and proposition quantifiers
provide further settings for testing that heuristic; ordinary algorithmic
decidability is not a question about arbitrary set-sized formulas.

\begin{hbdefinition}{def:p08-bvb501-abstract}{Finite occurrence, expansion, depth and relativization}
An abstract modal language has formulas with satisfaction at pointed
models. Each formula has a finite relevant signature; truth depends only
on that reduct and is unchanged by expanding other symbols. It extends
basic modal logic when every basic formula has an equivalent abstract
formula. Finite depth means that for some \(k\) depending on the
formula, truth at the point agrees with truth in the induced submodel
of points reachable in at most \(k\) steps. Bisimulation invariance
is an additional property, not a consequence of finite occurrence.

For the compactness characterization, require Boolean closure, including
negation, and relativization: \(\operatorname{Rel}(\chi,p)\)
is true iff \(p\) holds at the selected point and \(\chi\)
holds in the induced \(p\)-submodel. This qualification prevents
undefined evaluation when the point is outside that submodel. Compactness
means every finitely satisfiable formula set is satisfiable. Each formula
has finite atom and modality occurrence, but different formulas may
involve different finite signatures. Expansion invariance permits the
proof to work in the selected formula's finite relevant signature;
no fixed finite bound on the language's entire modality alphabet is imposed.
\end{hbdefinition}

\begin{hbtheorem}{thm:p08-bvb501-modal-lindstrom}{proved=annex:p08.src025501.lindstrom}{Source Theorems 50 and 52 with their logical closures explicit}
A finite-occurrence, expansion-invariant abstract language extending
basic modal logic is equivalent in expressive power to that language
iff it is bisimulation invariant and has finite depth. If it additionally
has Boolean closure and relativization, finite depth follows from
compactness and bisimulation invariance. Consequently compactness and
invariance characterize the basic modal language in this setting.
\end{hbtheorem}

The source attributes the finite-depth formulation to de Rijke and the
relativization proof to van Benthem. It recalls two classical
Lindstr\"om characterizations of first-order logic among its abstract
extensions: compactness together with the L\"owenheim--Skolem property,
and, in the alternative formulation, compactness together with
potential-isomorphism invariance (the Karp property). Each uses its
standard abstract-logic closure standing. The universal modality shows
why the finite-depth criterion does not describe every richer modal language.

\begin{hbtheorem}{thm:p08-bvb501-guarded-lindstrom}{cited=BlackburnVanBenthem2006ILLC}{Source Theorem 53: the guarded analog}
For the appropriate relativizing abstract guarded logics extending the
guarded fragment, equivalence in expressive power to the guarded fragment
is characterized by compactness and invariance under guarded bisimulation.
Satisfaction here is at guarded tuples, with the abstract-logic finite-
occurrence, expansion and Boolean closures; it is not just the preceding
one-point unary interface. The source refers the proof to van Benthem's
PP-2006-06. The guarded theorem is a separately cited result, not a
word-for-word consequence of the unary proof.
\end{hbtheorem}

\section{Alternative semantics}
\label{sec:p08-bvb501-alternatives}

\subsection{Boolean algebras with operators}
\label{sec:p08-bvb501-algebra}
\index{algebraic semantics}
\index{Boolean algebras with operators}
\index{BAOs}
\index{duality theory}
\index{discriminator term}

\begin{hbdefinition}{def:p08-bvb501-bao}{The source algebraic signature}
A Boolean algebra \((B,+,\times,-,1,0)\) is a bounded distributive
complemented lattice: \(+\) is join, \(\times\) meet, and
\(-\) complement. In particular \(b\times1=b\),
\(b+0=b\), \(b+(-b)=1\), and \(b\times(-b)=0\).
A Boolean algebra with operator has a unary \(m\) satisfying
\(m(b+c)=m(b)+m(c)\) and \(m(0)=0\). It interprets diamond;
box is \(-m(-b)\). A valuation sends atoms to elements and extends
by the corresponding algebraic operations. Validity means value \(1\)
under every valuation. Thus the join law mirrors
\(\Diamond(\varphi\lor\psi)\leftrightarrow
(\Diamond\varphi\lor\Diamond\psi)\), with both diamonds present.
\end{hbdefinition}

\begin{hbtheorem}{thm:p08-bvb501-algebra-completeness}{proved=annex:p08.src025501.algebra}{Source Theorem 55 and its normal-extension argument}
K is complete for all Boolean algebras with operators. Every normal
extension \(L\) is complete for the algebras validating its theorems.
Its Lindenbaum--Tarski algebra has classes of formulas modulo
\(L\)-provable equivalence and supplies a countervaluation for every
non-theorem. This does not imply relational-frame completeness of \(L\).
\end{hbtheorem}

The source relates representation theorems to completeness: Stone
representation gives Boolean algebras as fields of sets; J\'onsson--Tarski
representation embeds Boolean modal algebras into relational powerset
algebras with a distinguished field of admissible sets. Preservation is
for that field's valuations, not automatically for every subset valuation
of the resulting frame. This distinction permits algebraic completeness
without contradicting the relational incompleteness examples. Cayley's
permutation representation is the source's comparison from group theory.

For complex algebras of frames, disjoint union corresponds to a product
of algebras, generated-subframe restriction to a surjective homomorphism,
and a surjective bounded morphism to an injective inverse-image
homomorphism. Ultrafilter extensions connect the Boolean representation
to its powerset completion; their truth lemma was proved earlier.
The source explains the algebraic Goldblatt--Thomason argument through
Birkhoff's homomorphic-image/subalgebra/product characterization, and
cites Blok's splitting analysis of the lattice of normal modal logics.
It also notes that the universal diamond on a relational powerset algebra
sends the empty set to empty and every nonempty set to the whole domain,
a zero-discriminator operation. Algebraic methods can study non-Boolean
propositional bases too; Boolean structure is the present choice, not
a requirement on every modality studied algebraically.

\subsection{Neighborhood truth sets}
\label{sec:p08-bvb501-neighborhood}
\index{neighborhood model}
\index{neighborhood semantics}

\begin{hbdefinition}{def:p08-bvb501-neighborhood}{The source's point-to-subset relation}
A model has nonempty \(W\), valuation \(V\), and a relation
\(R\subseteq W\times\mathcal P(W)\); here \(R\) is not
a binary world accessibility relation. Put
\(N_w=\{X\subseteq W:wRX\}\). Then
\(M,w\entails\Box\chi\) iff its truth set belongs to \(N_w\),
and \(M,w\entails\Diamond\chi\) iff the complement of its
truth set does not belong to \(N_w\). No filter or upward-closure
assumption is imposed. A relational model embeds by the principal
successor filter \(N_w=\{X:R[w]\subseteq X\}\), not the singleton
family \(\{R[w]\}\).
\end{hbdefinition}

\begin{hbtheorem}{thm:p08-bvb501-neighborhood-results}{proved=annex:p08.src025501.neighborhood}{The corrected embedding, replacement law and source Theorem 56}
The principal-filter embedding preserves all basic modal truth. In
arbitrary neighborhoods K and theorem necessitation can fail, while
replacement of globally equivalent arguments preserves box equivalence:
explicitly, validity of \(\varphi\leftrightarrow\psi\) implies
validity of \(\Box\varphi\leftrightarrow\Box\psi\).
Neighborhood satisfiability is NP-complete and has a model with at most
\(1+r(r-1)/2\) points, where \(r\) is the number of distinct
box arguments after diamond duality expansion. This is a polynomial
bound in the input length. It is neighborhood, not unrestricted
relational-K satisfiability, that has this NP classification.
\end{hbtheorem}

In the source's Figure~17, \(W=\{t,u,v,w\}\),
\(V(p)=\{t,u\}\), \(V(q)=\{u,v\}\), and
\(N_u=\{\{t,u\},\{u,v,w\}\}\). Thus both \(\Box p\)
and \(\Box(p\to q)\) are true at \(u\), but \(\Box q\)
is false there. Other neighborhoods may be empty. On a singleton with
empty neighborhood family, \(\topv\) is true while \(\Box\topv\)
is false, disproving necessitation. The neighborhood membership framework
does not force logical omniscience. Upward closure, intersection closure
and the whole-domain condition specify stronger systems, as in Chellas's
distinct presentations. The source motivates weaker modal behavior by
strategic ability: forcing a disjunction need not mean being able to
force a selected disjunct. It credits Montague, Scott and Segerberg for
the neighborhood framework and Vardi for the complexity analysis.

\subsection{Topological interiors, closures and bisimulations}
\label{sec:p08-bvb501-topology}
\index{topological space}
\index{topological model}
\index{topological semantics}
\index{open set}
\index{open neighborhood}
\index{closed set}
\index{discrete topology}
\index{trivial topology}
\index{continuous maps}
\index{open maps}
\index{connected space}

\begin{hbdefinition}{def:p08-bvb501-topology}{Topological modal models and open matching}
A topology \(\tau\) on nonempty \(W\) contains \(W,\varnothing\),
is closed under finite intersections and arbitrary unions. A model
\((W,\tau,V)\) interprets box as interior of the argument truth
set and diamond as its closure. Equivalently box truth requires an
open neighborhood wholly satisfying the argument; diamond truth requires
an argument witness in every open neighborhood. Discrete topology is
\(\mathcal P(W)\), trivial topology \(\{\varnothing,W\}\).

A topo-bisimulation is a nonempty atomically harmonious relation such
that at a related pair \(wEw'\), each open neighborhood \(O\)
of \(w\) has an open neighborhood \(O'\) of \(w'\)
whose every point is related to some point of \(O\); symmetrically
every \(O'\) has an \(O\) whose every point has a related
point in \(O'\). These are the quantifiers of source Figure~18;
it does not depict only a single witness pair.
\end{hbdefinition}

\begin{hbtheorem}{thm:p08-bvb501-topological-invariance}{proved=annex:p08.src025501.topology}{Source Lemma 58, open restriction, maps and finite converse}
Topo-bisimulations preserve all basic modal formulas. Restriction to a
nonempty open subspace preserves truth at retained points. An open
continuous function gives such a bisimulation through its graph when
the source valuation is the inverse image of the target valuation;
arbitrary direct images of valuations do not suffice. For finite
topological models, modal equivalence at points implies topo-bisimilarity.
The interior interpretation validates T and 4, as well as K and
theorem necessitation.
\end{hbtheorem}

\begin{hbtheorem}{thm:p08-bvb501-topological-completeness}{cited=BlackburnVanBenthem2006ILLC}{The source's topological S4 results}
S4 is complete for topological spaces and has the finite topological
model property. The source also reports the McKinsey--Tarski result
that the real line, and every nonempty separable metric space without
isolated points, has precisely S4 as its basic modal validities.
It cites Esakia and Shehtman for further topological developments and
bases its account on Aiello, van Benthem and Bezhanishvili.
\end{hbtheorem}

\begin{hbtheorem}{thm:p08-bvb501-topological-definability}{proved=annex:p08.src025501.topology-def}{Discrete, connected and separation conditions}
Space validity of \(p\leftrightarrow\Box p\), quantified over
all valuations, defines discrete topology. Connectedness is not basic
modally definable, but is defined with the universal modality by
\(A(\Diamond p\to\Box p)\to(Ap\lor A\neg p)\).
The hybrid formula \(@_i\neg j\to(@_j\Box\neg i\lor @_i\Box\neg j)\)
defines T0; \(i\leftrightarrow\Diamond i\) defines T1.
\end{hbtheorem}

The source cites Gabelaia and Sustretov for the corresponding general
definability results, and T2 undefinability in the specified basic hybrid
language. A statement about one fixed valuation is not a topology's
all-valuation validity condition.

\section{The changing environment of modal reasoning}
\label{sec:p08-bvb501-environment}

The source's closing account connects philosophical modalities,
knowledge, conditionals and obligation with mathematical time, space
and provability; computational processes and verification; AI knowledge
representation, nonmonotonic reasoning and revision; linguistic structure;
and agency, security, games and social-choice powers. These are application
settings for language design, not evidence that one minimal language
expresses every task in them. Its internet/tree and intelligent-network
examples reflect its 2006 setting.

The unifying argument is mathematical rather than terminological:
model truth and frame validity have different correspondence languages;
bisimulations organize which distinctions a language retains;
translations expose first-order, second-order and infinitary resources;
and guarding or fixed-point discipline can retain useful computation.
Combination and first-order quantification show where properties cease
to transfer. Algebraic representation and topological bisimulation give
additional routes back to this structural account. The source's suggested
future emphasis on games and coalgebra is a research expectation, not
a theorem about modal expressivity.

Its acknowledgments credit Carlos Areces for editorial advice, Eric Kow
for diagrams and formatting, and Sara Hobin and Dmitry Sustretov for
the penultimate manuscript. The institutional manifestation, its
acknowledgments and its source pagination are separate from both a
revised course manuscript and the publisher-final chapter.

\section{Reading and topic map for the institutional source}
\label{sec:p08-bvb501-reading-map}

The numbered identities in this map refer specifically to the bibliography
of ILLC PP-2006-30, pp.~76--79 \cite{BlackburnVanBenthem2006ILLC},
not to the handbook's independently ordered bibliography. They identify
all 133 source references by their role in this presentation. A cited
external theorem remains a theorem attributed to its source; this
reading map does not identify different editions or papers as one work.

The foundational and historical chain comprises source references
31, 51, 60, 70--71, 75--78, 86--87, 94--96 and 133:
Dummett--Lemmon, Goldblatt's historical account, Hintikka, Kanger,
Kripke, Montague, Prior, Quine and von Wright. Their roles concern the
development and interpretation of modal reasoning, not a single shared
formal calculus.

Structural semantics, correspondence and abstract model theory draw on
5, 8--11, 13, 21, 27--29, 37--38, 52--55, 79, 82--83, 89--90,
98, 100, 113--114, 116--119, 121--122 and 124--127. These include
Andr\'eka--van Benthem--N\'emeti, Baltag, Barwise--Feferman,
Barwise--Moss, Barwise--van Benthem, Blackburn--de Rijke--Venema,
Chang--Keisler, de Rijke, Fine, Goldblatt--Thomason, Goranko--Passy,
Kurtonina--de Rijke, Lindstr\"om, Makinson, Otto, Park, Rosen,
Sahlqvist, Thomason and van Benthem's distinct correspondence,
ultrafilter, dynamics, definability and abstract-modal works.

Computation and representation use 3, 12, 16--17, 19--20, 23, 26,
30, 32--33, 35, 41--42, 56--59, 65--66, 72--73, 80, 88,
92--93, 107--108, 115 and 128--130: Akers, Berger, Bryant,
Burch--Clarke--McMillan--Dill--Hwang, Chagrova, Chandra--Merlin,
Clarke--Grumberg--Peled, De Giacomo--Lenzerini,
Dovier--Piazza--Policriti, Ebbinghaus--Flum, Emerson--Jutla,
Fagin--Halpern--Moses--Vardi, Fischer--Ladner, Gr\"adel,
Gr\"adel--Walukiewicz, Halpern--Vardi, Hemaspaandra, Huth--Ryan,
Janin--Walukiewicz, Kozen--Parikh, Kozen, Lee, Moore, Pratt,
Spaan, Streett--Emerson, Urquhart and Vardi. Algorithmic classifications
involving different inputs and semantic restrictions remain distinct.

Richer-language and combination routes use 4, 6--7, 14, 18, 24--25,
36, 39--40, 43--47, 49--50, 61--64, 69, 74, 81, 84, 91,
97, 99, 101--104, 109, 112, 120, 123 and 131--132. The chain
includes Andr\'eka--Hodkinson--N\'emeti,
Areces--Blackburn--Marx, Blackburn--Marx, Burgess,
Corsi--Ghilardi, Cresswell, Fine, Fine--Schurz, Fitting--Mendelsohn,
Friedman--Halpern, Gabbay and his temporal collaborators,
G\"ardenfors--Rott, Ghilardi, Hodkinson--Otto,
Hoogland--Marx--Otto, Kamp, Kracht--Wolter, Lewis, Marx,
Passy--Tinchev, Restall, Ryan--Schobbens, Scott, Segerberg,
Shehtman--Skvortsov, Sturm, ten Cate, van Benthem, Veltman and
Wolter. Quantified, guarded, conditional, hybrid and product settings
are not interchangeable versions of the basic language.

Alternative semantics and spatial definability use 1--2, 15, 22, 34,
48, 67--68, 85, 105--106 and 110--111: Aiello--van Benthem,
Aiello--van Benthem--Bezhanishvili, Blok, Chellas, Esakia,
Gabelaia, J\'onsson--Tarski's two papers, McKinsey--Tarski,
Shehtman and Sustretov's distinct proceedings/report manifestations.
These supplement the earlier neighborhood and representation references.

For navigation, the source's subject index on pp.~81--82 is organized
here by the actual argument locations rather than inherited page numbers:
\begin{center}
\begin{tabular}{p{0.62\linewidth}p{0.28\linewidth}}
Signature, worlds, valuations, truth, box/diamond, local/global consequence
 & Section~\ref{sec:p08-bvb501-basic}\\
Temporal, epistemic, extensional examples; isomorphism; evaluation games
 & Section~\ref{sec:p08-bvb501-graph-examples}\\
Standard translation, correspondence predicates, compactness, cardinal transfer
 & Section~\ref{sec:p08-bvb501-standard-translation}\\
Bisimulation, contraction, unraveling, generated models, bounded morphisms
 & Section~\ref{sec:p08-bvb501-bisimulation}\\
Structural comparison, infinitary definitions, interpolation and comparison games
 & Sections~\ref{sec:p08-bvb501-first-order-comparison}--\ref{sec:p08-bvb501-comparison-games}\\
Checking, labelling, two variables, filtration, finite models, depth and complexity
 & Section~\ref{sec:p08-bvb501-computation}\\
Frame validity, normal extensions, incompleteness, correspondents and ultrafilters
 & Section~\ref{sec:p08-bvb501-frame-logics}\\
Universal/difference, nominals, satisfaction operators, temporal and conditional logics
 & Sections~\ref{sec:p08-bvb501-universal}--\ref{sec:p08-bvb501-conditional}\\
Guarded/packed, programs, induction, fixed points, combinations and quantified modalities
 & Sections~\ref{sec:p08-bvb501-guarded}--\ref{sec:p08-bvb501-first-order-modal}\\
Abstract logics, guarding and characterization boundaries
 & Section~\ref{sec:p08-bvb501-abstract}\\
BAOs, representation, duality, neighborhoods, topological truth and separation
 & Section~\ref{sec:p08-bvb501-alternatives}
\end{tabular}
\end{center}
